% =====================================================================================
% Projektbericht "Webscraper für Modulhandbücher mit AWS"  (Methoden des Machine Learning, SS 2026)
% Bauen:  siehe bericht/README.md  (latexmk -pdf main.tex  bzw.  bericht/build.ps1)
% Zahlen und Grafiken kommen aus  zahlen.tex / bilder/  (erzeugt von make_report_charts.py).
% Offene Stellen sind mit \TODO{...} gelb markiert, Platzhalter für Sascha Lauk orange.
% =====================================================================================
\documentclass[a4paper,11pt,headings=small,parskip=half-,numbers=noenddot,bibliography=totoc]{scrreprt}

\usepackage[T1]{fontenc}
\usepackage[utf8]{inputenc}
\usepackage{lmodern}
\usepackage[ngerman]{babel}
\usepackage[autostyle]{csquotes}
\usepackage{microtype}
\usepackage[htt]{hyphenat}
\usepackage{xurl}
\usepackage[top=2.6cm,bottom=2.6cm,left=2.8cm,right=2.4cm]{geometry}
\usepackage{graphicx}
\usepackage{booktabs}
\usepackage{longtable}
\usepackage{tabularx}
\usepackage{array}
\usepackage{xcolor}
\usepackage{enumitem}
\usepackage{caption}
\usepackage{float}
\usepackage{rotating}
\usepackage{siunitx}
\usepackage{amsmath}
\usepackage{listings}
\usepackage[hang]{footmisc}
\usepackage{tikz}
\usetikzlibrary{positioning,fit,arrows.meta,shapes.geometric,backgrounds,calc}
\usepackage[numbers]{natbib}
\usepackage[hidelinks,pdfusetitle]{hyperref}
\usepackage[nameinlink,ngerman]{cleveref}

\renewcommand{\_}{\textunderscore\allowbreak{}}
\sisetup{locale=DE,group-separator={.},group-minimum-digits=4}
\captionsetup{font=small,labelfont=bf,format=hang}
\setlist{itemsep=1pt,topsep=3pt}
\setlength{\emergencystretch}{2em}
\renewcommand{\arraystretch}{1.15}
\graphicspath{{bilder/}}

\lstset{basicstyle=\ttfamily\footnotesize,breaklines=true,frame=single,
  framerule=0.3pt,rulecolor=\color{black!25},backgroundcolor=\color{black!3},
  columns=fullflexible,keepspaces=true,showstringspaces=false,
  literate={ä}{{\"a}}1 {ö}{{\"o}}1 {ü}{{\"u}}1 {Ä}{{\"A}}1 {Ö}{{\"O}}1 {Ü}{{\"U}}1 {ß}{{\ss}}1}

% ---- Markierungen für offene Stellen ------------------------------------------------
\newcommand{\TODO}[1]{\par\smallskip\noindent\colorbox{yellow!70}{\parbox{0.94\linewidth}{\footnotesize\textbf{TODO:} #1}}\par\smallskip}
\newcommand{\saschaPH}[1]{\par\smallskip\noindent\colorbox{orange!35}{\parbox{0.94\linewidth}{\footnotesize\textbf{[Sascha Lauk:} #1\textbf{]}}}\par\smallskip}
\newcommand{\saschaCell}{\colorbox{orange!35}{\footnotesize [Sascha]}}
\newcommand{\schaetzung}{\textsuperscript{\textit{(Schätzung)}}}

% ---- Zahlen & Diagramme ------------------------------------------------------------
% AUTOMATISCH ERZEUGT von bericht/make_report_charts.py - nicht von Hand editieren.
\newcommand{\zCovFH}{60}
\newcommand{\zCovKirch}{51}
\newcommand{\zCovKunst}{70}
\newcommand{\zCovOeff}{78}
\newcommand{\zCovPriv}{39}
\newcommand{\zCovUni}{76}
\newcommand{\zCovVerw}{44}
\newcommand{\zDauerAug}{6,3}
\newcommand{\zDauerDeep}{12,0}
\newcommand{\zDocsProUniGross}{909}
\newcommand{\zDocsProUniKlein}{147}
\newcommand{\zDocsTotal}{164.396}
\newcommand{\zDupLikely}{2.840}
\newcommand{\zDupStrict}{1.420}
\newcommand{\zDupVolumePct}{15}
\newcommand{\zEstTrue}{18.703}
\newcommand{\zEstTrueHigh}{19.172}
\newcommand{\zEstTrueLow}{18.150}
\newcommand{\zFPTotal}{17.972}
\newcommand{\zFinalListed}{21.108}
\newcommand{\zFinalStrict}{7.367}
\newcommand{\zGBHochrechnung}{280}
\newcommand{\zInitialPositive}{44.437}
\newcommand{\zJahrAnteilNeu}{61}
\newcommand{\zJahrAnzahl}{11.290}
\newcommand{\zJahrMedian}{2021}
\newcommand{\zJahrMin}{2000}
\newcommand{\zJahrMinAnzahl}{2}
\newcommand{\zJahreAlt}{2.037}
\newcommand{\zJahreMitte}{2.364}
\newcommand{\zKostenBedrock}{75}
\newcommand{\zKostenEc}{45}
\newcommand{\zKostenExtern}{0}
\newcommand{\zKostenFargate}{8}
\newcommand{\zKostenGesamt}{140}
\newcommand{\zKostenLambda}{2}
\newcommand{\zKostenSonstiges}{4}
\newcommand{\zKostenSpeicher}{6}
\newcommand{\zLLMConfirmed}{7.367}
\newcommand{\zLLMDocsHochrechnung}{59.195}
\newcommand{\zLLMKostenHochrechnung}{160}
\newcommand{\zLLMRejected}{20.308}
\newcommand{\zLLMReviewed}{27.678}
\newcommand{\zMHProUniGross}{235}
\newcommand{\zMHProUniGrossMax}{158}
\newcommand{\zMHProUniKlein}{7,1}
\newcommand{\zMHStrictVolumeGB}{6,3}
\newcommand{\zMHVolumeGB}{21,2}
\newcommand{\zManuellMitH}{15}
\newcommand{\zManuellOhneH}{231}
\newcommand{\zMeanMB}{0,80}
\newcommand{\zMedianMBAll}{0,24}
\newcommand{\zMedianMBMH}{0,46}
\newcommand{\zMedianMH}{4}
\newcommand{\zMidPrec}{25}
\newcommand{\zOCRConf}{42}
\newcommand{\zOCRDone}{3.553}
\newcommand{\zOCRLLM}{247}
\newcommand{\zOCRText}{3.118}
\newcommand{\zOCRTotal}{3.619}
\newcommand{\zPDFsHochrechnung}{351.593}
\newcommand{\zPDFsProMH}{8,8}
\newcommand{\zPilotHi}{85,9}
\newcommand{\zPilotLo}{78,5}
\newcommand{\zPilotN}{400}
\newcommand{\zPilotOK}{330}
\newcommand{\zPilotPct}{82,5}
\newcommand{\zRateAug}{10.467}
\newcommand{\zRateDeep}{11.016}
\newcommand{\zRecallMissed}{383}
\newcommand{\zSpeicherGB}{154,1}
\newcommand{\zSpeicherGiB}{143,5}
\newcommand{\zSpeicherObjekte}{193.438}
\newcommand{\zStundenJan}{156}
\newcommand{\zStundenPhaseDrei}{46}
\newcommand{\zStundenPhaseEins}{34}
\newcommand{\zStundenPhaseVier}{30}
\newcommand{\zStundenPhaseZwei}{46}
\newcommand{\zTfidfHigh}{13.741}
\newcommand{\zTfidfNegative}{114.926}
\newcommand{\zTfidfOther}{6.393}
\newcommand{\zTiefeMHBisVier}{65}
\newcommand{\zTiefeModusAlle}{3}
\newcommand{\zTiefeModusMH}{5}
\newcommand{\zTopTenShare}{37}
\newcommand{\zTrainPrec}{96,5}
\newcommand{\zUniqueGB}{131,0}
\newcommand{\zUniqueLikely}{18.268}
\newcommand{\zUniqueStrict}{5.947}
\newcommand{\zUnisGeTen}{161}
\newcommand{\zUnisLikely}{263}
\newcommand{\zUnisLikelyPct}{65}
\newcommand{\zUnisNone}{39}
\newcommand{\zUnisStrict}{248}
\newcommand{\zUnisStrictPct}{61}
\newcommand{\zUnisTotal}{404}
\newcommand{\zUnresolved}{1.658}
\newcommand{\zZielwert}{40.000}

% TikZ-Stile und Diagramme des Berichts: Datenfluss, Zeitleiste (Architektur: architektur.tex).
% Farbkonvention: grau = bestehend, blau = in dieser Version neu, orange = Sackgasse/ersetzt.

\tikzset{
  comp/.style   = {draw=black!55, rounded corners=2pt, fill=white, align=center, font=\footnotesize, inner sep=3pt, minimum height=0.75cm},
  newc/.style   = {comp, draw=blue!70!black, fill=blue!10, line width=0.8pt},
  oldc/.style   = {comp, draw=orange!80!black, fill=orange!10, dashed},
  store/.style  = {draw=black!55, fill=white, shape=cylinder, shape border rotate=90, aspect=0.25, align=center, font=\footnotesize, minimum height=0.95cm, minimum width=1.5cm},
  newstore/.style = {store, draw=blue!70!black, fill=blue!10, line width=0.8pt},
  ext/.style    = {comp, draw=black!40, fill=black!5, dotted},
  layer/.style  = {draw=black!35, rounded corners=4pt, inner sep=7pt},
  newlayer/.style = {layer, draw=blue!60!black},
  ll/.style     = {font=\scriptsize\itshape, text=black!65},
  arr/.style    = {-{Stealth[length=2.2mm]}, black!70, line width=0.5pt},
  narr/.style   = {-{Stealth[length=2.2mm]}, blue!70!black, line width=0.8pt},
  darr/.style   = {-{Stealth[length=2.2mm]}, black!55, dashed},
  lbl/.style    = {font=\tiny, text=black!70, fill=white, inner sep=1pt},
}

% ---------------------------------------------------------------- Datenfluss je Dokument ---
\newcommand{\datenfluss}{%
\begin{tikzpicture}[x=1cm,y=1cm,every node/.style={transform shape},scale=0.9]
  \node[comp,minimum width=2.2cm] (a) at (0,0) {Crawl\\\zDocsTotal\ PDFs};
  \node[comp,minimum width=2.2cm] (b) at (3.6,0) {TF-IDF-Modell\\Score 0--1};
  \node[comp,minimum width=2.2cm] (c) at (7.4,1.4) {Score $\ge$ 0,9\\\zTfidfHigh\ (nicht geprüft)};
  \node[comp,minimum width=2.2cm] (d) at (7.4,-0.2) {Score 0,5--0,9\\LLM-Review};
  \node[comp,minimum width=2.2cm] (e) at (7.4,-1.8) {Score $<$ 0,5\\\zTfidfNegative\ verworfen};
  \node[newc,minimum width=2.4cm] (f) at (12.2,-0.2) {\zLLMConfirmed\ bestätigt\\\zLLMRejected\ verworfen};
  \node[newc,minimum width=2.4cm] (g) at (12.2,-2.2) {OCR bei Scans\\\zOCRDone\ bearbeitet};
  \node[comp,minimum width=2.2cm,draw=blue!70!black,line width=1pt] (h) at (15.8,0.6) {Endliste\\\zFinalListed};
  \draw[arr] (a) -- (b);
  \draw[arr] (b) -- (c);
  \draw[arr] (b) -- (d);
  \draw[arr] (b) -- (e);
  \draw[narr] (d) -- (f);
  \draw[narr] (e) -- node[lbl,below]{leere Textebene} (g);
  \draw[narr] (g) -- (f);
  \draw[arr] (c) -| (h);
  \draw[narr] (f) -- (h);
\end{tikzpicture}}

% ---------------------------------------------------------------- Zeitleiste ---------------
\newcommand{\zlw}{0.9}
\newcommand{\zlms}[3]{\draw[black!60] (#1*\zlw,0.55) -- (#1*\zlw,#3); \node[anchor=south,align=center,font=\scriptsize,inner sep=1pt] at (#1*\zlw,#3) {#2};}
\newcommand{\zeitleiste}{%
\begin{tikzpicture}[x=1cm,y=1cm,every node/.style={font=\scriptsize}]
  \foreach \k [count=\i from 0] in {24,...,40} {
    \draw[black!30] (\i*\zlw,0) -- (\i*\zlw,-0.12);
    \node[black!60,anchor=north] at (\i*\zlw+0.45*\zlw,-0.1) {\k};
  }
  \draw[black!50] (0,0) -- (17*\zlw,0);
  \node[black!60,anchor=east] at (-0.1,-0.35) {KW};
  \fill[blue!60,rounded corners=1pt] (0*\zlw,0.15) rectangle (4*\zlw,0.55);
  \fill[green!55!black!70,rounded corners=1pt] (4*\zlw,0.15) rectangle (7*\zlw,0.55);
  \fill[yellow!70!orange,rounded corners=1pt] (8*\zlw,0.15) rectangle (13*\zlw,0.55);
  \fill[violet!70,rounded corners=1pt] (14*\zlw,0.15) rectangle (17*\zlw,0.55);
  \node[white] at (2*\zlw,0.35) {Phase 1};
  \node[white] at (5.5*\zlw,0.35) {Phase 2};
  \node[black] at (10.5*\zlw,0.35) {Phase 3};
  \node[white] at (15.5*\zlw,0.35) {Phase 4};
  \zlms{0.7}{10.06.\\Projektstart,\\Scrapy-Prototyp}{0.85}
  \zlms{2.6}{24.06.\\Lambda, S3,\\DynamoDB}{2.0}
  \zlms{3.5}{03.07.\\EC2-Webapp,\\Authentifizierung}{3.15}
  \zlms{5.0}{12.07.\\Klassifikator,\\Review-UI}{0.85}
  \zlms{6.0}{22.07.\\CI/CD,\\Feedback-Loop}{2.0}
  \zlms{6.8}{24.07.\\Fargate,\\Best-First-Crawl}{3.15}
  \zlms{9.5}{12.08.\\\textbf{Sandbox verloren}\\Neuaufbau, Großlauf}{0.85}
  \zlms{11.6}{27.08.\\Serper, OCR,\\LLM-Review}{2.0}
  \zlms{12.6}{03.09.\\Seeds-only\\(+43 Hochschulen)}{3.15}
  \zlms{14.9}{16.09.\\\textbf{Budget-Sperre}\\Deep Run (12\,h)}{0.85}
  \zlms{16.4}{30.09.\\LLM-Review,\\Auswertung}{2.0}
\end{tikzpicture}}

% Architekturdiagramme v1-v3 (TikZ). Styles stehen in diagramme.tex.

% ---------------------------------------------------------------- Architektur v1 -----------
\newcommand{\archA}{%
\begin{tikzpicture}[x=1cm,y=1cm,every node/.style={transform shape},scale=0.92]
  \node[ext] (user) at (0,0) {Nutzer\\Browser};
  \node[comp] (caddy) at (2.5,0) {Caddy\\HTTPS};
  \node[comp] (api) at (5.2,0.6) {FastAPI\\+ uvicorn};
  \node[store] (sqlite) at (5.2,-1.6) {SQLite\\Nutzer/Rollen};
  \begin{scope}[on background layer]
    \node[layer,fit=(caddy)(api)(sqlite),label={[ll]above:EC2 t3.small (Docker Compose)}] (ec2) {};
  \end{scope}
  \node[comp] (jr) at (9.8,0.6) {JobRunner\\(1 Job je URL)};
  \node[comp] (spider) at (12.6,0.6) {DocumentSpider\\Scrapy, Tiefe 2};
  \node[comp] (pipe) at (15.4,0.6) {Pipelines\\Local/S3};
  \begin{scope}[on background layer]
    \node[layer,fit=(jr)(spider)(pipe),label={[ll]above:AWS Lambda (Container-Image, 15\,min Limit)}] (lam) {};
  \end{scope}
  \node[store] (ddb) at (9.8,-2.3) {DynamoDB\\Visited-Store};
  \node[ext] (web) at (12.6,-2.3) {Hochschul-\\webseiten};
  \node[store] (s3) at (15.4,-2.3) {S3\\PDFs};
  \draw[arr] (user) -- (caddy);
  \draw[arr] (caddy) -- (api);
  \draw[arr] (api) -- (sqlite);
  \draw[arr] (api) -- node[lbl,above]{invoke (synchron)} (jr);
  \draw[arr] (jr) -- (spider);
  \draw[arr] (spider) -- (pipe);
  \draw[arr] (pipe) -- (s3);
  \draw[arr] (jr) -- (ddb);
  \draw[arr] (spider) -- (web);
\end{tikzpicture}}

% ---------------------------------------------------------------- Architektur v2 -----------
\newcommand{\archB}{%
\begin{tikzpicture}[x=1cm,y=1cm,every node/.style={transform shape},scale=0.86]
  \node[newc] (train) at (0,4.8) {Entwickler-Laptop\\Training \texttt{mlclassifier}};
  \node[newc] (gh) at (4.6,4.8) {GitHub Actions\\(OIDC)};
  \node[newc] (ecr) at (8.2,4.8) {ECR\\Container-Image};
  \node[newc] (ssm) at (4.6,3.2) {SSM\\Run Command};
  \node[ext] (user) at (0,0.8) {Nutzer\\Browser};
  \node[comp] (caddy) at (2.3,0.8) {Caddy};
  \node[comp] (api) at (4.6,1.5) {FastAPI};
  \node[newc] (rev) at (4.6,-0.3) {Review-UI\\Feedback};
  \node[store] (sqlite) at (2.3,-1.1) {SQLite};
  \begin{scope}[on background layer]
    \node[layer,fit=(caddy)(api)(rev)(sqlite),label={[ll]below:EC2 (Webapp)}] (ec2) {};
  \end{scope}
  \node[newc] (route) at (7.6,0.8) {Routing nach\\URL-Anzahl};
  \node[comp] (lam) at (10.6,2.2) {Lambda\\$\le$10 URLs};
  \node[newc] (far) at (10.6,-0.2) {Fargate\\Bulk-Crawl};
  \node[newc] (cls) at (13.8,1.0) {Klassifikation\\TF-IDF + LogReg};
  \begin{scope}[on background layer]
    \node[layer,fit=(lam)(far)(cls),label={[ll]below:gleiches Container-Image}] (cmp) {};
  \end{scope}
  \node[newstore] (fb) at (4.6,-3.4) {S3 feedback/\\Verdikte};
  \node[store] (ddb) at (9.0,-3.4) {DynamoDB\\Visited};
  \node[store] (s3) at (13.8,-3.4) {S3\\PDFs, Manifeste};
  \draw[narr] (train) -- node[lbl,above]{git push (Modell)} (gh);
  \draw[narr] (gh) -- (ecr);
  \draw[narr] (gh) -- (ssm);
  \draw[narr] (ssm) -- (ssm |- ec2.north);
  \draw[narr] (ecr) -- node[lbl,right,pos=0.45]{update-function-code} (lam.north);
  \draw[arr] (user) -- (caddy);
  \draw[arr] (caddy) -- (api);
  \draw[arr] (api) -- (route);
  \draw[narr] (route) -- (lam);
  \draw[narr] (route) -- (far);
  \draw[arr] (lam) -- (cls);
  \draw[narr] (far) -- (cls);
  \draw[arr] (cls) -- (s3);
  \draw[narr] (rev) -- (fb);
  \draw[narr] (fb.west) -- ++(-5.4,0) |- node[lbl,pos=0.25,left]{Nachtraining} (train.west);
  \draw[arr] (far) -- (ddb);
\end{tikzpicture}}

% ---------------------------------------------------------------- Architektur v3 -----------
\newcommand{\archC}{%
\begin{tikzpicture}[x=1cm,y=1cm,every node/.style={transform shape},scale=0.80]
  \node[newc] (cc) at (0,5.2) {Common Crawl\\CDX-Index};
  \node[newc] (serp) at (0,3.9) {Serper\\(Google-SERP)};
  \node[newc] (fire) at (0,2.6) {Firecrawl\\(JS-Rendering)};
  \node[oldc] (gcse) at (0,1.2) {Google CSE,\\DuckDuckGo};
  \node[newc] (disc) at (3.4,4.4) {Discovery-Modul\\Seeds};
  \node[newstore] (seeds) at (7.0,4.7) {S3 discovery/\\Seed-Listen};
  \node[ext] (web) at (3.4,0.2) {Hochschul-\\webseiten};
  \node[newc] (spider) at (7.2,0.9) {Spider: Best-First,\\Sitemap, Sektions-Cap};
  \node[newc] (seedmode) at (10.4,2.4) {Seeds-only-Modus\\Robots-Bypass};
  \node[comp] (cls) at (10.4,0.6) {Klassifikation\\TF-IDF};
  \begin{scope}[on background layer]
    \node[layer,fit=(spider)(seedmode)(cls),label={[ll]below:Fargate-Task (4 vCPU, 16\,GB, Watchdog)}] (far) {};
  \end{scope}
  \node[store] (ddb) at (7.2,-2.2) {DynamoDB\\Visited};
  \node[store] (s3) at (13.8,0.8) {S3\\193.438 Objekte\\154\,GB};
  \node[newc] (ocr) at (13.8,4.0) {OCR (Laptop)\\Tesseract};
  \node[newc] (llm) at (13.8,-1.9) {LLM-Review (Laptop)\\Bedrock Haiku 4.5};
  \node[newc] (eval) at (10.4,-3.7) {Evaluation \&\\Auswertung (23 Grafiken)};
  \node[comp] (webapp) at (3.4,-2.4) {Webapp (EC2)\\Schalter: Discovery,\\Rendering};
  \node[newc] (boot) at (3.4,-4.6) {Bootstrap-Skript\\(baut alles neu auf)};
  \draw[narr] (cc) -- (disc);
  \draw[narr] (serp) -- (disc);
  \draw[narr] (fire) -- (disc);
  \draw[darr,orange!80!black] (gcse) -- (disc);
  \draw[narr] (disc) -- (seeds);
  \draw[narr] (seeds) -- (seedmode);
  \draw[arr] (spider) -- (cls);
  \draw[arr] (seedmode) -- (cls);
  \draw[arr] (web) -- (spider);
  \draw[arr] (spider) -- (ddb);
  \draw[arr] (cls) -- (s3);
  \draw[narr] (s3) -- (ocr);
  \draw[narr] (s3) -- (llm);
  \draw[narr] (llm) -- (eval);
  \draw[narr] (ocr.east) to[out=0,in=0,looseness=1.3] (llm.east);
  \draw[arr] (webapp.north) -- (web.south);
  \draw[narr] (boot) -- (webapp);
\end{tikzpicture}}


% Hilfsmakro: Abbildung mit PNG
\newcommand{\abb}[4][0.95]{%
  \begin{figure}[#2]\centering
    \includegraphics[width=#1\linewidth]{#3}
    \caption{#4}\label{fig:#3}
  \end{figure}}

\title{Projektbericht\\[0.4em]\Large Webscraper für Modulhandbücher mit AWS\\[0.2em]\large Skalierbare Erhebung und Klassifikation von Modulhandbüchern deutscher Hochschulen}
\author{Jan Soetebeer \and Sascha Lauk}
\date{Oktober 2026}
\subject{Modul „Methoden des Machine Learning“ \textperiodcentered{} FH Wedel \textperiodcentered{} Sommersemester 2026}

\begin{document}

\begin{titlepage}
  \centering
  \vspace*{1.5cm}
  {\large FH Wedel \textperiodcentered{} Master Informatik\par}
  {\normalsize Modul „Methoden des Machine Learning“ \textperiodcentered{} Sommersemester 2026\par}
  \vspace{2.2cm}
  {\Huge\bfseries Projektbericht\par}
  \vspace{0.8cm}
  {\LARGE Webscraper für Modulhandbücher mit AWS\par}
  \vspace{0.5cm}
  {\large Skalierbare Erhebung und Klassifikation von\\Modulhandbüchern deutscher Hochschulen\par}
  \vspace{2.5cm}
  {\large Jan Soetebeer \quad\textperiodcentered\quad Sascha Lauk\par}
  \vspace{1.2cm}
  {\normalsize Oktober 2026\par}
  \vfill
  {\small Entstanden im Kontext der Forschungsgruppe KIWIT (BMFTR-gefördert), Forschungsprofil A (Makroebene)\par}
  \vspace{0.3cm}
  {\footnotesize Entwurfsstand: Zahlen, die noch von laufenden LLM-Prüfungen abhängen, sind aus \texttt{auswertung/data/kennzahlen.json} generiert\\und werden vor der Abgabe aktualisiert (siehe Anhang~\ref{app:repro}).\par}
\end{titlepage}

\pagenumbering{roman}

\chapter*{Zusammenfassung}
\addcontentsline{toc}{chapter}{Zusammenfassung}
Dieser Bericht dokumentiert Entwicklung, Betrieb und Evaluation eines modularen Web-Crawlers, der öffentlich zugängliche Modulhandbücher deutscher Hochschulen erhebt und mit Methoden des Machine Learning von anderen PDF-Dokumenten unterscheidet. Der Bericht beschreibt nicht nur das Ergebnis, sondern vor allem den Weg dorthin: drei Architekturgenerationen auf AWS (von einer einzelnen Lambda-Funktion über Fargate-Bulk-Läufe bis zu einer Pipeline mit Discovery, OCR und LLM-Zweitstufe), die eingeschlagenen und wieder verworfenen Ansätze sowie die Rückschläge im Betrieb, darunter der Verlust der AWS-Sandbox und die Budget-Sperre.

Im Endstand wurden \zDocsTotal{} verschiedene PDFs von \zUnisTotal{} Hochschulen heruntergeladen (\zSpeicherObjekte{} S3-Objekte, \zSpeicherGB{}\,GB; zwei Großläufe von \zDauerAug{}\,h und \zDauerDeep{}\,h). Der TF-IDF-Klassifikator erreichte in der gruppierten Kreuzvalidierung eine Precision von \zTrainPrec{}\,\%, im realen Crawl jedoch nur \zMidPrec{}\,\% (Score 0,5--0,9) bzw. \zPilotPct{}\,\% (Score $\ge$ 0,9). Erst eine LLM-Zweitstufe (Claude Haiku 4.5 auf Amazon Bedrock) für \zLLMReviewed{} Dokumente brachte belastbare Qualität: \zLLMConfirmed{} Modulhandbücher sind LLM-bestätigt, zusammen mit den ungeprüften Treffern oberhalb von Score 0,9 ergibt sich eine Endliste von \zFinalListed{} Dokumenten, von denen schätzungsweise \zEstTrue{} echte Modulhandbücher sind. \zUnisLikely{} von \zUnisTotal{} Hochschulen (\zUnisLikelyPct{}\,\%) sind abgedeckt. Der Bericht diskutiert Grenzen (Recall, LLM als Referenz, Kosten) und gibt eine Hochrechnung für eine deutschlandweite Vollerhebung.

\vspace{0.8em}
\noindent\textbf{Schlüsselwörter:} Web Scraping, Scrapy, AWS Lambda, AWS Fargate, Dokumentklassifikation, TF-IDF, LLM-Review, Amazon Bedrock, Modulhandbücher.

\tableofcontents
\listoffigures
\listoftables

\clearpage
\pagenumbering{arabic}

\chapter{Einleitung}\label{ch:einleitung}

\section{Forschungskontext: KIWIT}
Das vorliegende Projekt ist im Kontext der Forschungsgruppe KIWIT (\emph{Funktionen und Folgen Künstlicher Intelligenz in der Wissenschafts- und Hochschulorganisation}) entstanden, die seit April 2025 vom Bundesministerium für Forschung, Technologie und Raumfahrt (BMFTR) gefördert wird. Die Beschreibung des Kontexts stammt vom Betreuer des Projekts:

\begin{quote}\small
„Das vorliegende Projekt ist im Kontext der Forschungsgruppe KIWIT (Funktionen und Folgen Künstlicher Intelligenz in der Wissenschafts- und Hochschulorganisation) entstanden, die seit April 2025 vom Bundesministerium für Forschung, Technologie und Raumfahrt (BMFTR) gefördert wird. KIWIT untersucht empirisch, wie Künstliche Intelligenz Strukturen und Prozesse in Wissenschaft und Hochschulorganisationen verändert --- auf der Ebene der Forschungspraxis, der institutionellen Steuerung sowie des Lehrens und Lernens.

Innerhalb dieses Rahmens ist der entwickelte Crawler dem Forschungsprofil A (Makroebene) zugeordnet und dient der systematischen Erhebung öffentlich zugänglicher Hochschulinformationen. Universitätswebseiten sind eine zentrale primäre Datenquelle, um zu erfassen, wie Hochschulen KI-bezogene Strategien, Richtlinien und Angebote kommunizieren und formalisieren. Das modular aufgebaute Tool ermöglicht eine skalierbare, reproduzierbare Datenerhebung über eine Vielzahl von Institutionen hinweg und bildet damit eine methodische Grundlage für die quantitativen Analysevorhaben innerhalb von KIWIT.“
\end{quote}

Daraus ergibt sich eine Anforderung, die den gesamten Projektverlauf geprägt hat: Das System soll \emph{keine Einmallösung} für genau eine Dokumentart sein, sondern eine \emph{wiederverwendbare Erhebungsplattform}. Modulhandbücher dienen als erster, gut überprüfbarer Anwendungsfall. Sie eignen sich dafür, weil sie auf den Webseiten fast aller Hochschulen als PDF vorliegen, aber nirgends zentral verzeichnet sind, und weil sich ihre Echtheit vergleichsweise eindeutig beurteilen lässt.

\section{Ausgangssituation und Motivation}
An der FH Wedel werden regelmäßig neue Studienangebote konzipiert. Für eine fundierte Entscheidungsgrundlage ist der Vergleich mit Studiengängen anderer Hochschulen unerlässlich, da deren Modulhandbücher verlässlich Auskunft über Lehrinhalte, Umfänge und Prüfungsformen geben. Bislang mussten solche Dokumente manuell auf den Webseiten der Hochschulen gesucht und gesichtet werden. Das ist zeitaufwändig und fehleranfällig: In einer Aufwandsabschätzung zum MVP wurden 10 bis 15 Minuten je Hochschule veranschlagt, für einen Durchlauf über alle Hochschulen also 40 bis 60 Stunden, und das Ergebnis müsste in jedem Semester neu erhoben werden.

Die Problemstellung ist dabei weniger das Herunterladen von Dokumenten, das Scrapy seit Jahren beherrscht, als das \emph{Finden} und \emph{Aussortieren}: Hochschulwebseiten sind sehr unterschiedlich aufgebaut, Handbücher liegen auf Fakultäts-Subdomains, in Modul-Datenbanken oder hinter skriptgenerierten Download-Links, und Prüfungsordnungen, Flyer oder Vorlesungsverzeichnisse sehen Modulhandbüchern täuschend ähnlich.

\section{Problemstellung und Zielsetzung}
Als User Story formuliert: \emph{Als Mitarbeiter der FH Wedel möchte ich einen Webscraper, der übergebene Links automatisiert nach Modulhandbüchern durchsucht, relevante von irrelevanten Dokumenten unterscheidet und die gefundenen Modulhandbücher strukturiert speichert, damit ich Lehrinhalte und -umfänge unterschiedlicher Hochschulen im Vergleich zu unseren Studienangeboten analysieren kann.}

Daraus wurden zu Projektbeginn vier Akzeptanzkriterien abgeleitet:
\begin{enumerate}[label=\textbf{A\arabic*}]
  \item Der Scraper durchsucht alle übergebenen Links vollständig nach Dokumenten, die als Modulhandbuch identifizierbar sind.
  \item Alle tatsächlichen Modulhandbücher werden gefunden --- keine relevanten Dokumente werden übersehen (Recall).
  \item Dokumente, die keine Modulhandbücher sind, werden zuverlässig erkannt und nicht gespeichert (Precision).
  \item Die identifizierten Modulhandbücher werden strukturiert und für die Weiterverarbeitung nutzbar abgelegt.
\end{enumerate}

Der Projektumfang hat sich im Verlauf deutlich verschoben. Der MVP im Juli arbeitete mit rund 6.400 PDFs von etwa 239 Hochschulen und einem Klassifikator, der auf 202 Dokumenten von 21 Hochschulen trainiert war. Ab August lautete das Ziel, \emph{möglichst alle} Modulhandbücher \emph{aller} \zUnisTotal{} Hochschulen der Hochschulliste zu erheben. Als grobe Fachschätzung für die Größenordnung diente ein Erwartungswert von etwa \zZielwert{} Modulhandbüchern. Dieser Wert ist keine Messung, sondern nur Referenz für die Recall-Einordnung.

\section{Leitfragen und Kennzahlen}
Neben der Funktion soll der Bericht den Daten- und Forschungsaspekt des Projekts greifbar machen. Dazu werden folgende Fragen mit Zahlen beantwortet (Kapitel~\ref{ch:ergebnisse}):
\begin{itemize}
  \item Wie lange dauert ein Scraping-Lauf, und wie viele Dokumente und Daten fallen dabei an?
  \item Wie viele Hochschulen werden erfasst, und wie unterscheiden sich kleine und große, öffentliche und private Hochschulen?
  \item Wie hoch ist die Erkennungs- bzw. Klassifikationsquote, und wie viele Dokumente müssen manuell nachbearbeitet werden?
  \item Wie viele Duplikate treten auf, wie alt sind die ältesten Dokumente, und in welcher Tiefe der Webseiten liegen die meisten Funde?
  \item Welche Datenmenge, Laufzeit und welche Kosten wären bei einer deutschlandweiten Vollerhebung zu erwarten?
\end{itemize}

\section{Vorgehen und Aufbau des Berichts}
Der Bericht folgt dem tatsächlichen Projektverlauf und stellt \emph{den Weg} gleichberechtigt neben das Ergebnis. Kapitel~\ref{ch:verlauf} gibt den zeitlichen Überblick mit wochenweiser Aufschlüsselung der Aktivitäten. Kapitel~\ref{ch:grundlagen} fasst die eingesetzten Technologien zusammen. Kapitel~\ref{ch:architektur} zeigt die Architektur in drei Versionen, Kapitel~\ref{ch:aws} die Entwicklung der AWS-Umgebung einschließlich der Probleme (Sandbox-Verlust, Budget-Sperre). Kapitel~\ref{ch:crawler} und~\ref{ch:klassifikation} beschreiben Crawler und Klassifikation mit allen eingeschlagenen Wegen. Kapitel~\ref{ch:ergebnisse} stellt die Kennzahlen dar, Kapitel~\ref{ch:diskussion} diskutiert Sackgassen, Lehren und Grenzen. Kapitel~\ref{ch:sascha} ist dem Beitrag von Sascha Lauk vorbehalten, Kapitel~\ref{ch:fazit} schließt mit Fazit und Ausblick.

\paragraph{Lesehinweis zu Farben in Diagrammen.} In den Architekturdiagrammen sind \emph{neu hinzugekommene} Komponenten blau hinterlegt, verworfene Ansätze orange gestrichelt. In den Ergebnisgrafiken gilt ein festes Schema: blau = LLM-bestätigt, grün = TF-IDF $\ge$ 0,9 (nicht LLM-geprüft), orange = verworfen, grau = negativ.

\chapter{Projektverlauf}\label{ch:verlauf}

Das Projekt lief von der ersten Zeile Code am 10.\,Juni 2026 bis zur Abschlussauswertung am 2.\,Oktober 2026, also über 17 Kalenderwochen. Rückblickend lässt es sich in vier Phasen gliedern, die sich jeweils durch ein dominierendes Problem auszeichnen (\cref{fig:zeitleiste}).

\begin{figure}[htbp]\centering
  \resizebox{\linewidth}{!}{\zeitleiste}
  \caption{Zeitleiste des Projekts mit Phasen und Meilensteinen (Kalenderwochen 24--40, 2026). Die Phasen sind in den folgenden Abschnitten beschrieben.}\label{fig:zeitleiste}
\end{figure}

\section{Phasen im Überblick}

\paragraph{Phase 1: Prototyp und AWS-Grundlage (KW\,24--27, Juni bis Anfang Juli).}
Ein Scrapy-Spider, der PDF- und Word-Dokumente von einer Start-URL herunterlädt, wird um einen \texttt{JobRunner} (ein Job je URL, parallel begrenzt) und einen persistenten Visited-Store erweitert. Das erste Architekturdiagramm entsteht (Architektur v1, \cref{sec:arch-v1}). Danach folgt die Cloud-Anbindung: Lambda-Container-Image in ECR, S3 für Dokumente, DynamoDB für den Visited-Store, anschließend Webapp mit Authentifizierung auf einer EC2-Instanz. Dominierendes Thema: \emph{überhaupt etwas Lauffähiges in die Cloud bringen}.

\paragraph{Phase 2: Klassifikator und Plattform (KW\,28--30, Juli).}
Der Endlosschleifen-Bug bei Mehrfach-Formaten wird behoben, Jobs laufen asynchron mit Fortschrittsanzeige, ein Terraform-Stack soll die Infrastruktur beschreiben. Der erste TF-IDF-Klassifikator wird trainiert und in die Scraping-Pipeline integriert. Es folgen die Human-in-the-Loop-Feedbackschleife, CI/CD über GitHub Actions und, als Antwort auf das 15-Minuten-Limit von Lambda, der Fargate-Bulk-Pfad (Architektur v2, \cref{sec:arch-v2}). Dominierendes Thema: \emph{Qualität und Skalierung des Crawls}.

\paragraph{Phase 3: Großlauf und Recall-Analyse (KW\,32--36, August bis Anfang September).}
Am 12.\,August läuft die AWS-Sandbox ab, alle Daten sind verloren; die Umgebung wird über ein Bootstrap-Skript neu aufgebaut, und noch am selben Tag wird der Großlauf über 411 Hochschulen in der neuen Umgebung wiederholt. Die Auswertung zeigt, dass nur ein Teil der erwarteten Handbücher gefunden wurde. Es entstehen das Evaluationsmodul, die Discovery (Common Crawl, Serper, Firecrawl), der OCR-Fallback, die LLM-Zweitstufe und der Seeds-only-Modus (Architektur v3, \cref{sec:arch-v3}). Dominierendes Thema: \emph{Recall}.

\paragraph{Phase 4: Deep Run und LLM-Bereinigung (KW\,38--40, September bis Oktober).}
Nach der Budget-Sperre der Sandbox wird der Crawl um Breitensuche und Laufzeit-Wachhund erweitert und als zwölfstündiger Deep Run gestartet (\zDocsTotal{} PDFs kumuliert). Weil die TF-IDF-Positive in großer Zahl falsch sind, werden rund \zLLMReviewed{} Dokumente von einem LLM geprüft, Scans per OCR nachbearbeitet und die Abschlussauswertung mit 23 Grafiken erstellt. Dominierendes Thema: \emph{Precision}.

\section{Wochenweise Aktivität}
\Cref{tab:wochen-jan} schlüsselt die Aktivität von Jan Soetebeer wochenweise auf. Die Aufgaben stammen aus der Commit-Historie (28 Commits auf \texttt{main}, 42 inklusive Nebenzweigen) und den Entwicklungssitzungen. Die Stundenangaben sind \emph{Schätzungen}; sie wurden anhand der dokumentierten Sitzungs- und Commit-Zeiten sowie des Aufwands für Warte- und Betriebszeiten (Läufe, AWS-Konsole, Recherche) angesetzt \TODO{Stunden von Jan prüfen und anpassen (Quelle: bericht/daten/stunden.csv)}.

{\small
\begin{longtable}{@{}l l p{8.9cm} r@{}}
\caption{Wochenübersicht Jan Soetebeer (Stunden geschätzt).}\label{tab:wochen-jan}\\
\toprule
\textbf{Woche} & \textbf{Zeitraum} & \textbf{Tätigkeiten} & \textbf{Std.} \\
\midrule
\endfirsthead
\toprule
\textbf{Woche} & \textbf{Zeitraum} & \textbf{Tätigkeiten} & \textbf{Std.} \\
\midrule
\endhead
\bottomrule
\endlastfoot
% AUTOMATISCH ERZEUGT aus daten/stunden.csv
KW\,24 & 08.06.--14.06. & Projektstart; Scrapy-Prototyp (DocumentSpider, PDF/Word-Download); JobRunner mit parallelen Jobs; Visited-Store; erstes Architekturdiagramm (Architektur v1) & ca.\,10 \\
KW\,25 & 15.06.--21.06. & Einarbeitung AWS (Lambda, ECR, IAM); keine Commits [Jan: prüfen] & ca.\,2 \\
KW\,26 & 22.06.--28.06. & AWS-Integration: Lambda-Container-Image, ECR, S3-Pipeline, DynamoDB-Visited-Store, IAM-Policies; Frontend-Grundgerüst (Login, Tabs) & ca.\,12 \\
KW\,27 & 29.06.--05.07. & Authentifizierung (bcrypt, Session-Token); Hosting der Webapp auf EC2 mit Docker Compose und Caddy; Rollen- und Modellverwaltung & ca.\,10 \\
KW\,28 & 06.07.--12.07. & Endlosschleifen-Bug (Timeout bei Mehrfach-Formaten) behoben; asynchrone Jobs mit Fortschrittsanzeige; Terraform-Stack; Tiefen-Crawl; Ergebnis- und Review-UI; erster Klassifikator (TF-IDF + logistische Regression, Gruppen-Kreuzvalidierung); Scraper-Integration & ca.\,16 \\
KW\,29 & 13.07.--19.07. & MVP-Präsentation erstellt (15.07.-16.07.); Deployment-Versuche Lambda-Image und EC2; Klassifikator in Lambda-Image eingebettet & ca.\,8 \\
KW\,30 & 20.07.--26.07. & Human-in-the-Loop-Feedbackschleife (Review-UI, S3-Feedback, Nachtraining); CI/CD mit GitHub Actions (OIDC, SSM); CSV-Import der Hochschulliste; Skalierungs-Roadmap; fokussierter Crawl (Best-First, Sitemap, Subdomains); Extraktionsprofile; Fargate-Bulk-Pfad (Architektur v2) & ca.\,22 \\
KW\,31 & 27.07.--02.08. & keine dokumentierte Aktivität & ca.\,0 \\
KW\,32 & 03.08.--09.08. & Planung der Evaluation des ersten Großlaufs & ca.\,3 \\
KW\,33 & 10.08.--16.08. & Ablauf der Sandbox: kompletter Neuaufbau der AWS-Umgebung (Bootstrap-Skript); erster Großlauf über 411 Hochschulen (6,3 h); Evaluationsmodul mit HTML-Bericht & ca.\,14 \\
KW\,34 & 17.08.--23.08. & Recall-Analyse (145 Hochschulen ohne Treffer); Schwellenwert-Anpassung; Discovery-Modul (Common Crawl); OCR-Fallback & ca.\,8 \\
KW\,35 & 24.08.--30.08. & Serper- und Firecrawl-Anbindung; LLM-Zweitstufe auf Amazon Bedrock; End-to-End-Test auf Beispielhochschulen & ca.\,9 \\
KW\,36 & 31.08.--06.09. & Webapp-Schalter für Discovery/Rendering; parallele Bedrock-Aufrufe; Nachtraining mit 5.000+ Modulhandbüchern; Seeds-only-Modus; Per-Hochschule-Diagnose; Seeds-only-Lauf (+43 Hochschulen) & ca.\,12 \\
KW\,37 & 07.09.--13.09. & keine dokumentierte Aktivität (abwesend) & ca.\,0 \\
KW\,38 & 14.09.--20.09. & Budget-Sperre der Sandbox analysiert; Robots-Bypass für Discovery-Seeds; Fehler in der Evaluation (LLM-Urteile überschrieben) behoben; Breadth-Fairness und Watchdog; Deep Run (12 h, 132.193 PDFs) & ca.\,14 \\
KW\,39 & 21.09.--27.09. & keine dokumentierte Aktivität & ca.\,0 \\
KW\,40 & 28.09.--04.10. & LLM-Review von rund 27.700 Dokumenten; OCR von 3.553 Scans; Pilotstichprobe; Abschlussauswertung mit 23 Grafiken; Projektbericht & ca.\,16 \\
\midrule
\multicolumn{3}{l}{\textbf{Summe (Schätzung)}} & \textbf{ca.\,156} \\

\end{longtable}}

\Cref{fig:B01_stunden_pro_woche} zeigt dieselben Werte als Säulendiagramm. Die Spitzen liegen in KW\,30 (CI/CD, Feedback-Loop, fokussierter Crawl und Fargate an zwei Tagen), KW\,28 (Klassifikator und Review-UI) sowie in KW\,33 (Sandbox-Neuaufbau und erster Großlauf). Die Pausen in KW\,31, 37 und 39 sind Urlaubs- bzw. Abwesenheitszeiten, in denen nichts gelaufen ist; das ist für die Budget-Sperre in Kapitel~\ref{ch:aws} relevant. Insgesamt ergeben sich rund \zStundenJan{}~Stunden (Phase~1: \zStundenPhaseEins{}\,h, Phase~2: \zStundenPhaseZwei{}\,h, Phase~3: \zStundenPhaseDrei{}\,h, Phase~4: \zStundenPhaseVier{}\,h).

\abb[0.92]{htbp}{B01_stunden_pro_woche}{Geschätzter Arbeitsaufwand von Jan Soetebeer je Kalenderwoche, eingefärbt nach Projektphase.}

\subsection*{Wochenübersicht Sascha Lauk}
\saschaPH{Wochenweise Aktivität und Stunden ergänzen. Die Tabelle ist vorbereitet; Zeitraum und Wochen sind identisch zu Tabelle~\ref{tab:wochen-jan}.}

{\small
\begin{longtable}{@{}l l p{8.9cm} r@{}}
\caption{Wochenübersicht Sascha Lauk (Platzhalter).}\label{tab:wochen-sascha}\\
\toprule
\textbf{Woche} & \textbf{Zeitraum} & \textbf{Tätigkeiten} & \textbf{Std.} \\
\midrule
\endfirsthead
\toprule
\textbf{Woche} & \textbf{Zeitraum} & \textbf{Tätigkeiten} & \textbf{Std.} \\
\midrule
\endhead
\bottomrule
\endlastfoot
% AUTOMATISCH ERZEUGT aus daten/stunden.csv - Platzhalter für Sascha Lauk
KW\,24 & 08.06.--14.06. & \saschaCell{} & \saschaCell{} \\
KW\,25 & 15.06.--21.06. & \saschaCell{} & \saschaCell{} \\
KW\,26 & 22.06.--28.06. & \saschaCell{} & \saschaCell{} \\
KW\,27 & 29.06.--05.07. & \saschaCell{} & \saschaCell{} \\
KW\,28 & 06.07.--12.07. & \saschaCell{} & \saschaCell{} \\
KW\,29 & 13.07.--19.07. & \saschaCell{} & \saschaCell{} \\
KW\,30 & 20.07.--26.07. & \saschaCell{} & \saschaCell{} \\
KW\,31 & 27.07.--02.08. & \saschaCell{} & \saschaCell{} \\
KW\,32 & 03.08.--09.08. & \saschaCell{} & \saschaCell{} \\
KW\,33 & 10.08.--16.08. & \saschaCell{} & \saschaCell{} \\
KW\,34 & 17.08.--23.08. & \saschaCell{} & \saschaCell{} \\
KW\,35 & 24.08.--30.08. & \saschaCell{} & \saschaCell{} \\
KW\,36 & 31.08.--06.09. & \saschaCell{} & \saschaCell{} \\
KW\,37 & 07.09.--13.09. & \saschaCell{} & \saschaCell{} \\
KW\,38 & 14.09.--20.09. & \saschaCell{} & \saschaCell{} \\
KW\,39 & 21.09.--27.09. & \saschaCell{} & \saschaCell{} \\
KW\,40 & 28.09.--04.10. & \saschaCell{} & \saschaCell{} \\
\midrule
\multicolumn{3}{l}{\textbf{Summe}} & \saschaCell{} \\

\end{longtable}}

\section{Arbeitsteilung}
Das Projekt wurde von zwei Personen bearbeitet. Der vorliegende Hauptteil des Berichts beschreibt die Arbeiten von Jan Soetebeer: Scraper und Crawl-Strategie, AWS-Umgebung und Betrieb, Klassifikator, Discovery, OCR und LLM-Zweitstufe sowie die Abschlussauswertung. Der Beitrag von Sascha Lauk wird in Kapitel~\ref{ch:sascha} von ihm selbst dargestellt \saschaPH{Aufgabenbereich in zwei bis drei Sätzen}.

\section{Werkzeuge und Arbeitsweise}
Entwickelt wurde in Python~3.12 mit Git/GitHub (Branches \texttt{main}, \texttt{develop}, \texttt{feature/add-ML-Models}, Pull-Requests), Docker und der AWS-CLI. Bei der Implementierung wurde ein KI-Coding-Assistent (Claude Code) eingesetzt; die Sitzungsprotokolle dienten zusätzlich als Quelle für die Rekonstruktion des Projektverlaufs. In der LLM-Zweitstufe der Pipeline selbst kommt Claude Haiku~4.5 über Amazon Bedrock zum Einsatz (Abschnitt~\ref{sec:llm}). \TODO{Kennzeichnung des KI-Einsatzes mit den Vorgaben der FH Wedel abgleichen}

\paragraph{Umfang der Implementierung.} Das Repository umfasst rund 7.800 Zeilen Python im Crawler (\texttt{webscraper/}, 51 Dateien), rund 2.300 Zeilen im Klassifikator-Paket (\texttt{mlclassifier/}), rund 800 Zeilen in den Auswertungsskripten, rund 1.900 Zeilen Frontend-Code (HTML/CSS/JavaScript) und rund 1.800 Zeilen Infrastruktur- und Workflow-Dateien (Bootstrap-Skript, GitHub-Actions-Workflows, Terraform).

\chapter{Grundlagen und Technologiestack}\label{ch:grundlagen}

Leitend für alle Technologieentscheidungen war das Anforderungsprofil eines projektbasierten, aber auf Wiederverwendbarkeit ausgelegten Systems: geringe Betriebskosten, wenig Infrastrukturaufwand, kurze Entwicklungszeit und eine Architektur, die sich sauber erweitern lässt. \Cref{tab:stack} gibt den Gesamtüberblick; die Begründungen der wichtigsten Entscheidungen folgen in den Abschnitten.

\begin{table}[htbp]
\centering\small
\caption{Technologiestack im Überblick (Stand Projektende).}\label{tab:stack}
\begin{tabularx}{\linewidth}{@{}l X l@{}}
\toprule
\textbf{Bereich} & \textbf{Technologie und Rolle} & \textbf{Seit} \\
\midrule
Crawler & Scrapy (Spider, Pipelines, Dupe-Filter, AutoThrottle); generische Engine mit austauschbaren Extraktionsprofilen & Juni \\
Webapp & FastAPI + uvicorn, statisches HTML/CSS/JavaScript, SQLite (Nutzer, Rollen, Modelle), bcrypt, Caddy als HTTPS-Proxy & Juni/Juli \\
Rechnen & AWS Lambda (interaktiv, $\le$ 10 URLs) und AWS Fargate (Bulk); gleiches Container-Image & Juni/Juli \\
Zustand & S3 (Dokumente, Manifeste, Feedback, Seeds), DynamoDB (Visited-Store) & Juni \\
Klassifikation & scikit-learn: TF-IDF (Wort- und Zeichen-n-Gramme) + logistische Regression, PyMuPDF zur Textextraktion & Juli \\
Zweitstufe & Claude Haiku 4.5 über Amazon Bedrock (spezifikationsgesteuerter Prompt), Tesseract-OCR für Scans & August \\
Discovery & Common Crawl (CDX-Index), Serper (Google-SERP-API), Firecrawl (JS-Rendering) & August \\
Betrieb & GitHub Actions (OIDC, SSM), Bootstrap-Skript (PowerShell + AWS-CLI), Docker Compose; Terraform nur zeitweise & Juli/August \\
Auswertung & Python (matplotlib, stdlib-Evaluationspaket), HTML-Berichte, \LaTeX{} & August--Oktober \\
\bottomrule
\end{tabularx}
\end{table}

\section{Crawler: Scrapy}
Scrapy liefert die Bausteine für robustes, wiederholtes Crawlen: einen Dupe-Filter innerhalb eines Laufs, konfigurierbare Drosselung (AutoThrottle) und eine Pipeline-Architektur, in die sich Verarbeitungsschritte als austauschbare Klassen einhängen lassen. Genau diese Erweiterbarkeit war ausschlaggebend: Ein neuer Spider, eine neue Pipeline-Stufe oder ein neues Visited-Store-Backend sind jeweils eine zusätzliche Klasse. Die Höflichkeitsregeln (robots.txt nach RFC~9309 \citep{rfc9309}, Verzögerung, identifizierbarer User-Agent) sind Standard, nicht Option.

\section{Webapp: FastAPI und SQLite}
FastAPI liefert REST-API und statisches Frontend aus demselben Prozess; es entfällt eine separate Frontend-Infrastruktur. Router trennen die Domänen (\texttt{auth}, \texttt{users}, \texttt{roles}, \texttt{models}, \texttt{scrape}). Als Datenhaltung genügt SQLite auf einem Docker-Volume, weil nur ein Prozess schreibt und die Nutzerzahl einstellig ist; das Schema wird idempotent beim Start angelegt. Passwörter liegen als bcrypt-Hash vor, Sitzungen sind signierte 12-Stunden-Tokens. Das Rollenmodell trennt die \emph{Position} (Administrator/User) von fachlichen \emph{Rollen} (zunächst nur „Modulhandbuch“), und Modelle werden über dieselben Rollen Nutzern zugeordnet. Damit lässt sich ein neuer Anwendungsfall mit einer neuen Rolle und einem passenden Modell abbilden, ohne Frontend-Code zu ändern.

\section{Klassifikation: TF-IDF statt Transformer}
Für „Ist dieses PDF ein Modulhandbuch?“ wurde bewusst eine klassische Pipeline gewählt: TF-IDF \citep{salton1988} über Wort- und Zeichen-n-Gramme und logistische Regression \citep{pedregosa2011}. Zu Projektbeginn standen 202 Dokumente aus 21 Hochschulen zur Verfügung; dafür ist das Finetuning eines Transformers die falsche Werkzeugklasse, während ein regularisiertes lineares Modell zuverlässig funktioniert, in Sekunden trainiert ist und ohne GPU in Lambda und Fargate läuft. Der eigentliche Flaschenhals ist Netzwerk-I/O und PDF-Textextraktion, nicht Matrixmultiplikation. Dass diese Entscheidung für die Massenfilterung richtig, für die Endqualität aber nicht ausreichend war, zeigt Kapitel~\ref{ch:klassifikation}.

\subsection{Evaluationsmetriken}
Bei ungleichgewichtigen Klassen ist die Fläche unter der Precision-Recall-Kurve (PR-AUC, hier als \emph{average precision}) die Leitmetrik. Precision $P=\frac{TP}{TP+FP}$ gibt an, welcher Anteil der als Modulhandbuch gemeldeten Dokumente wirklich eines ist; Recall $R=\frac{TP}{TP+FN}$, welcher Anteil der echten gefunden wurde. Unsicherheiten von Anteilen werden als Wilson-Intervall \citep{wilson1927} angegeben. Die Schwellen werden nicht auf 0,5 gesetzt, sondern aus den Daten auf ein Recall-Ziel von 95\,\% kalibriert, weil ein übersehenes Handbuch eine Datenlücke ist, ein Fehlalarm dagegen nur wenige Sekunden Review kostet. Die Entscheidung ist dreistufig: \texttt{automatic\_negative}, \texttt{needs\_review} und \texttt{automatic\_positive}.

\section{Cloud-Komponenten}
Die Wahl von AWS ergab sich aus dem bereitgestellten Hochschul-Account (AWS-Sandbox der FH Wedel). Warum zunächst Lambda gewählt wurde und warum es später durch Fargate ergänzt werden musste, ist Gegenstand von Kapitel~\ref{ch:aws}. Wichtig für die Plattformidee: \emph{ein} Container-Image (Crawler + \texttt{mlclassifier} + Modell) läuft in allen Laufzeitumgebungen (Lambda, Fargate, lokal), nur der Einstiegspunkt unterscheidet sich.

\section{Externe Dienste}
\begin{itemize}
  \item \textbf{Common Crawl} \citep{commoncrawl}: freier Webindex; über den CDX-Index lassen sich pro Domain PDF-URLs mit „modul/handbuch“-Mustern finden, ohne die Hochschul-Server zu belasten.
  \item \textbf{Serper}: Google-SERP-API für Abfragen der Form \texttt{site:<domain> <Begriff>}; erreicht tief verlinkte Dateien, die Common Crawl nicht kennt. Es konnte nur die kostenlose Stufe genutzt werden.
  \item \textbf{Firecrawl}: rendert JavaScript-Seiten und liefert Link-Karten; im Projekt vorbereitet, aber nur ansatzweise getestet.
  \item \textbf{Amazon Bedrock} mit Claude Haiku 4.5 \citep{anthropic2025haiku}: Zweitstufen-Klassifikation für die Grauzone.
  \item \textbf{Tesseract}: OCR für gescannte PDFs ohne Textebene.
\end{itemize}

\chapter{Architektur und ihre Entwicklung}\label{ch:architektur}

Die Architektur ist nicht in einem Wurf entstanden, sondern in drei Generationen gewachsen, die jeweils auf ein konkretes Problem der vorherigen reagiert haben. Dieses Kapitel zeigt jede Version als Diagramm (neue Komponenten blau, verworfene orange) und begründet die Übergänge. Die zugehörigen AWS-Details (Dienste, IAM, Kosten, Betriebsprobleme) stehen in Kapitel~\ref{ch:aws}.

\section{Leitprinzipien}
Über alle Versionen hinweg gelten vier Prinzipien:
\begin{enumerate}
  \item \textbf{Steuerung und Last trennen.} Das „Cockpit“ (Webapp auf EC2) ist klein, läuft dauerhaft und skaliert nicht mit der Datenmenge. Die Arbeit läuft nur bei Bedarf in kurzlebigen Containern.
  \item \textbf{Ein Image, mehrere Laufzeiten.} Derselbe Container-Code läuft in Lambda, Fargate und lokal; der Einstiegspunkt (Lambda-Handler bzw. \texttt{bulk\_run.py}) unterscheidet sich.
  \item \textbf{Generische Engine, austauschbare Profile.} Die Crawl-Engine kennt keine Modulhandbücher; sie delegiert vier Entscheidungen (Linkbewertung, Zieldefinition, Extraktion aus Binärdateien, Extraktion aus Seiten) an ein \emph{Extraktionsprofil}. Ein neuer Anwendungsfall ist ein neues Profil plus ein Registry-Eintrag.
  \item \textbf{Zustand extern halten.} Dokumente, Manifeste, Feedback, Seeds und Visited-Store liegen in S3 bzw. DynamoDB, nicht im Container.
\end{enumerate}

\section{Architektur v1: Lambda plus Webapp (Juni bis Anfang Juli)}\label{sec:arch-v1}
Die erste Version ist ein klassischer Serverless-Ansatz (\cref{fig:archA}). Die Webapp auf einer EC2-Instanz (FastAPI, SQLite, Caddy für HTTPS) nimmt eine URL oder eine CSV-Liste entgegen und ruft pro Anfrage eine Lambda-Funktion synchron auf. Die Funktion startet den \texttt{JobRunner}, der pro URL einen Scrapy-Spider betreibt (Tiefe 2, höchstens 60 Seiten), die Dateien über die Pipelines nach S3 schreibt und die besuchte URL im DynamoDB-Visited-Store vermerkt.

\begin{sidewaysfigure}\centering
  \resizebox{0.95\linewidth}{!}{\archA}
  \caption{Architektur v1 (Juni bis 3.\,Juli): Webapp auf EC2 ruft Lambda synchron auf; S3 und DynamoDB halten den Zustand.}\label{fig:archA}
\end{sidewaysfigure}

\paragraph{Begründung.} Die Last ist selten, stoßweise und perfekt parallelisierbar: 424 Hochschulen haben nichts miteinander zu tun. Lambda skaliert auf null und kostet ohne Last nichts. Rechnerisch ergaben sich für einen Vollauf 424 parallele Invocations (innerhalb des Limits von 1.000), ein Rechenaufwand von rund 76.000 GB-Sekunden und damit etwa 1,30~USD. Im MVP-Vortrag wurde ein Vollauf auf rund zehn Minuten veranschlagt; diese Zahl war \emph{gerechnet, nicht gemessen} (zur späteren Realität siehe Abschnitt~\ref{sec:aws-lambda-fargate}).

\paragraph{Grenzen.} Drei harte Lambda-Eigenschaften erwiesen sich als Problem: das Zeitlimit von 15 Minuten (Tiefe/Seitenzahl mussten klein bleiben, sonst Timeout), das Puffern jeder Datei im Arbeitsspeicher und die Tatsache, dass alle URLs eines Aufrufs sich ein Budget teilen. Die synchrone Antwort der Webapp lief bei längeren Läufen zudem zuverlässig in den HTTP-Timeout; deshalb wurde das Polling-Muster (POST liefert eine \texttt{job\_id}, GET fragt den Status) eingeführt.

\section{Architektur v2: Klassifikation, Feedback, CI/CD und Fargate (Juli)}\label{sec:arch-v2}
Die zweite Version (\cref{fig:archB}) fügt vier Dinge hinzu:
\begin{enumerate}
  \item \textbf{Klassifikation im Crawl.} Eine \texttt{ClassificationPipeline} bewertet jedes heruntergeladene PDF im Speicher mit dem TF-IDF-Modell, annotiert das Item und schreibt ein Manifest (eine JSONL-Zeile je Dokument mit Score und Entscheidung) nach S3. Das Modell ist in das Container-Image eingebettet; Training und Betrieb nutzen dasselbe Paket (\texttt{mlclassifier}), ein Train/Serve-Skew ist damit ausgeschlossen.
  \item \textbf{Human-in-the-Loop.} Die Review-UI zeigt Positive und Zweifelsfälle; Verdikte („Modulhandbuch“/„Kein Modulhandbuch“) werden gebündelt nach \texttt{s3://…/feedback/} geschrieben. Ein lokales Kommando (\texttt{feedback-retrain}) zieht die Verdikte, holt die Dokumente und trainiert nach.
  \item \textbf{CI/CD.} Ein Push auf \texttt{main} baut das Image und aktualisiert Lambda (GitHub Actions, OIDC statt statischer Schlüssel) und rollt die Webapp über SSM Run Command aus, ohne dass ein SSH-Port geöffnet wird.
  \item \textbf{Fargate-Bulk-Pfad.} Die Webapp routet nach URL-Anzahl: bis \texttt{BULK\_URL\_THRESHOLD} (zunächst 10) an Lambda, darüber an einen ECS-Fargate-Task mit demselben Image und dem Einstiegspunkt \texttt{bulk\_run.py}. Der Task hat kein Zeitlimit, mehr Speicher und liest die URL-Liste aus S3.
\end{enumerate}

\begin{sidewaysfigure}\centering
  \resizebox{0.95\linewidth}{!}{\archB}
  \caption{Architektur v2 (11.--24.\,Juli): Klassifikation im Container-Image, Feedback-Schleife, CI/CD und Fargate-Bulk-Pfad mit Routing nach URL-Anzahl.}\label{fig:archB}
\end{sidewaysfigure}

Im selben Zeitraum wurde die Crawl-Engine modularisiert (Best-First-Frontier, Sitemap-Seeding, Subdomain-Erkennung, Extraktionsprofile; Kapitel~\ref{ch:crawler}) und das Terraform-Projekt angelegt, das später verworfen wurde (Abschnitt~\ref{sec:aws-iac}).

\section{Architektur v3: Discovery, OCR und LLM-Zweitstufe (August bis Oktober)}\label{sec:arch-v3}
Die Auswertung des ersten Großlaufs (Kapitel~\ref{ch:crawler}) zeigte, dass \emph{Finden} das Hauptproblem ist, später, dass \emph{Aussortieren} es ist. Die dritte Version (\cref{fig:archC}) antwortet auf beides:
\begin{itemize}
  \item \textbf{Discovery-Modul.} Common Crawl und Serper liefern vorab URL-Kandidaten (\emph{Seeds}) pro Hochschule; sie liegen in S3 und werden vom Fargate-Task als priorisierte Ziele geladen. Der \emph{Seeds-only-Modus} holt ausschließlich diese Seeds und überspringt den Visited-Store, sodass nur Lücken-Hochschulen ohne Duplikate nachgearbeitet werden. Für Seeds lässt sich robots.txt gezielt umgehen (Abschnitt~\ref{sec:crawler-robots}).
  \item \textbf{Crawl-Verbesserungen.} Breadth-Fairness (Seitenkappe je Fakultäts-Sektion), größere Tiefe und Seitenzahl sowie ein Laufzeit-Wachhund, der den Task hart beendet.
  \item \textbf{Nachbereitung auf dem Laptop.} OCR (Tesseract) für gescannte PDFs, LLM-Review über Amazon Bedrock und die Auswertungsskripte arbeiten direkt auf S3-Daten.
  \item \textbf{Bootstrap-Skript.} Ein idempotentes PowerShell-Skript baut die gesamte AWS-Umgebung neu auf und ersetzt Terraform; es entstand als Reaktion auf den Sandbox-Verlust (Kapitel~\ref{ch:aws}).
  \item \textbf{Verworfene Quellen} (orange): Google Custom Search und DuckDuckGo (Abschnitt~\ref{sec:crawler-discovery}).
\end{itemize}

\begin{sidewaysfigure}\centering
  \resizebox{0.95\linewidth}{!}{\archC}
  \caption{Architektur v3 (August bis Oktober): Discovery-Seeds, Seeds-only-Modus, Fargate-Crawl mit Wachhund, OCR und LLM-Review auf S3-Daten sowie Bootstrap-Skript.}\label{fig:archC}
\end{sidewaysfigure}

\section{Vergleich der Versionen}
\Cref{tab:arch-vergleich} stellt die drei Versionen gegenüber. Auffällig ist, wie sich der Engpass verschoben hat: von der \emph{Plattform} (läuft es in der Cloud?) über die \emph{Laufzeit} (reicht Lambda?) und den \emph{Recall} (wird alles gefunden?) hin zur \emph{Precision} (ist es wirklich ein Modulhandbuch?).

\begin{table}[htbp]
\centering\small
\caption{Die drei Architekturversionen im Vergleich.}\label{tab:arch-vergleich}
\begin{tabularx}{\linewidth}{@{}p{2.6cm} X X X@{}}
\toprule
 & \textbf{v1 (Juni--Juli)} & \textbf{v2 (Juli)} & \textbf{v3 (August--Oktober)} \\
\midrule
Rechnen & Lambda (synchron) & Lambda + Fargate, Routing nach URL-Anzahl & Fargate-Bulk, Seeds-only; Nachbereitung lokal \\
Crawl & Tiefe 2, 60 Seiten, ungesteuert & Best-First, Sitemap, Subdomains, Profile & Seeds, Sektions-Cap, Tiefe 5, 2.500 Seiten \\
Klassifikation & -- & TF-IDF im Image, Review-Band & + LLM-Zweitstufe, OCR \\
Qualitätsschleife & -- & Human-in-the-Loop, Nachtraining & LLM-Labels, Pilotstichprobe \\
Deployment & manuell, Terraform-Entwurf & GitHub Actions (OIDC, SSM) & + Bootstrap-Skript für Neuaufbau \\
Haupt\-problem danach & 15 min, Speicher, wenig Funde & Recall (nur 261 von 404 Hochschulen) & Precision, Kosten, Betrieb \\
\bottomrule
\end{tabularx}
\end{table}

\section{Datenfluss pro Dokument}
Unabhängig von der Version durchläuft jedes Dokument dieselbe Kette (\cref{fig:datenfluss}): Download, Textextraktion, Bewertung durch das TF-IDF-Modell und je nach Score entweder direkte Übernahme, LLM-Prüfung oder Verwurf. Scans ohne Textebene gehen zunächst in die OCR und anschließend in dieselbe LLM-Prüfung. Die Zahlen sind die der Abschlussauswertung.

\begin{figure}[htbp]\centering
  \resizebox{\linewidth}{!}{\datenfluss}
  \caption{Datenfluss pro Dokument mit den Mengen der Abschlussauswertung (Endliste = LLM-bestätigt + TF-IDF $\ge$ 0,9 ungeprüft).}\label{fig:datenfluss}
\end{figure}

\section{Modularität und Wiederverwendbarkeit}
Die Forderung des Projektkontexts (Kapitel~\ref{ch:einleitung}), eine Plattform statt einer Einzellösung zu liefern, wurde an vier Stellen umgesetzt:
\begin{itemize}
  \item \textbf{Extraktionsprofile} (\texttt{profiles/}): \texttt{modulhandbuch} (Schlüsselwort-gewichtet, Referenzprofil), \texttt{html\_content} (Vorlage für HTML-Inhalte und strukturierte Felder). Der Engine-Code bleibt unverändert.
  \item \textbf{Spezifikationsgesteuerter LLM-Review} (\texttt{specs.py}): Eine \texttt{ClassificationSpec} (Definition, Beispiele, Ausschlüsse) beschreibt, was gesucht wird; das Urteilsfeld ist generisch (\texttt{is\_match}).
  \item \textbf{Profilgetriebene Discovery:} Suchbegriffe und URL-Stämme stammen aus dem Profil; ein neuer Dokumenttyp braucht nur ein anderes Profil.
  \item \textbf{Rollen-Modell-Matrix} der Webapp: neue Rolle plus neues Modell ergibt einen neuen Anwendungsfall ohne Frontend-Änderung.
\end{itemize}
Das ist die technische Brücke zu KIWIT: Sollen künftig statt Modulhandbüchern etwa KI-Strategiepapiere oder -Richtlinien erhoben werden, sind ein neues Profil, eine neue Spezifikation und ein neu trainiertes Modell nötig, nicht aber ein neuer Crawler.

\chapter{Die AWS-Umgebung: Aufbau, Entwicklung und Rückschläge}\label{ch:aws}

Die AWS-Umgebung war nicht nur Hosting, sondern ein zentraler Teil der Projektarbeit. Dieses Kapitel beschreibt, wie sich Rechenumgebung und Betrieb entwickelt haben (von Lambda zu Fargate, von Terraform zu einem Bootstrap-Skript), welche Beschränkungen des Hochschul-Accounts den Weg bestimmt haben und welche Rückschläge aufgetreten sind.

\section{Rahmenbedingungen: eine Sandbox auf Zeit}
Alle Ressourcen liefen in einer AWS-Sandbox der FH Wedel (Region \texttt{eu-central-1}, Zugang per SSO-Rolle). Drei Eigenschaften dieser Umgebung haben das Projekt wesentlich geprägt:
\begin{itemize}
  \item \textbf{Befristete Lease:} Läuft die Sandbox ab, wird \emph{alles} gelöscht. Die Konto-ID ändert sich (im Projekt von \texttt{660941536751} zu \texttt{081757578883}).
  \item \textbf{Budgetdeckel:} Für den Sandbox-Pool galt ein Limit von 50~USD; bei Überschreitung wird der Pool eingefroren.
  \item \textbf{Service Control Policies (SCP):} Einzelne Dienste und Aktionen sind organisationsweit gesperrt (Abschnitt~\ref{sec:aws-restriktionen}).
\end{itemize}

\section{Von Lambda zu Fargate (und warum nicht EC2 oder ECS)}\label{sec:aws-lambda-fargate}

\subsection{Ausgangsentscheidung}
Die ursprüngliche Begründung für Lambda war schlüssig: Scrape-Jobs sind selten, stoßweise und perfekt parallel, Lambda skaliert auf null, ein Dauerbetrieb über ECS oder EC2 würde „Abschalten bei Nichtnutzung“ erst nachbauen müssen. Im MVP-Vortrag wurde ECS/Fargate ausdrücklich als Lösung für ein Problem bezeichnet, das man nicht habe. Das stimmte für den Ausgangszustand (Jobs auf zehn Minuten gedeckelt, einige Dutzend URLs), nicht mehr für das Zielbild „alle Hochschulen, tief gecrawlt“.

\subsection{Wo Lambda scheiterte}
Beim Versuch, Tiefe und Seitenzahl zu erhöhen (Tiefe 1--3, 25--200 Seiten wurden durchprobiert), liefen die Jobs ins 15-Minuten-Limit. Am 22.\,Juli entstand deshalb eine Skalierungs-Roadmap. Sie identifizierte drei harte Grenzen von Lambda: (1)~das Zeitlimit von 15 Minuten, (2)~die Pufferung jeder Datei im Arbeitsspeicher (\texttt{upload\_fileobj(BytesIO(\ldots))}) und (3)~die gemeinsame Zeit- und Speicherbudgetierung aller URLs eines Aufrufs. Visited-Store (DynamoDB) und Ausgabe (S3) skalierten dagegen problemlos; das Nadelöhr war ausschließlich die Crawl-Rechenzeit.

\subsection{Optionen}
\Cref{tab:compute-optionen} fasst die erwogenen Rechenpfade zusammen.

\begin{table}[htbp]
\centering\small
\caption{Erwogene Rechenpfade für den Crawl und ihr Schicksal.}\label{tab:compute-optionen}
\begin{tabularx}{\linewidth}{@{}p{3.1cm} X p{3.0cm}@{}}
\toprule
\textbf{Option} & \textbf{Bewertung} & \textbf{Ergebnis} \\
\midrule
Lambda (nur) & Billig, schnell, aber 15\,min und RAM-Pufferung; ungeeignet für tiefe Crawls & behalten für $\le$ 10 URLs (interaktiv) \\
SQS-Fan-out auf Lambda & Eine Hochschule je Invocation, hunderte parallel; löst „viele Hochschulen“, nicht „eine sehr tiefe Hochschule“ & geplant (Roadmap, Phase~1b), \emph{nicht gebaut} \\
EC2-Stopgap & Großer Lauf auf der vorhandenen Webapp-Instanz; keine neue Infrastruktur, aber t3.small und belegt Dauerbetrieb & nur als Übergangslösung beschrieben \\
ECS auf EC2 & Cluster-Verwaltung nötig, Leerlaufkosten & nicht gewählt \\
\textbf{ECS Fargate} & Kein Zeitlimit, bis 4\,vCPU/16\,GB, gleiches Image, Abrechnung nur während der Laufzeit & \textbf{gewählt} (24.\,Juli) \\
\bottomrule
\end{tabularx}
\end{table}

\subsection{Umsetzung des Fargate-Pfads}
Am 24.\,Juli entstand der Bulk-Pfad innerhalb eines Tages: Einstiegspunkt \texttt{bulk\_run.py} (umweltvariablengesteuert), Task-Definition und IAM-Rollen als Vorlagen mit Platzhaltern, GitHub-Workflow \texttt{bulk-crawl.yml} (\texttt{workflow\_dispatch}) und ein URL-Routing in der Webapp. Die Defaults wurden schrittweise erhöht (2 vCPU/8\,GB~$\to$ 4\,vCPU/16\,GB; 300~$\to$ 800 Seiten; 15~$\to$ 30 Subdomain-Sitemaps). Hinzu kamen später ein Seeds-Modus und ein Wachhund (Abschnitt~\ref{sec:aws-kosten}). Eine Stolperfalle: Die Webapp scheiterte zunächst mit „TaskDefinition not found“, weil nur der Workflow die Task-Definition selbst registrierte; das Bootstrap-Skript registriert sie nun immer.

\subsection{Gerechnet versus gemessen}
\Cref{tab:rechnung-realitaet} stellt die Annahmen des MVP-Vortrags den tatsächlichen Läufen gegenüber. Die Lambda-Rechnung (10 Minuten für alle Hochschulen) beschrieb einen \emph{flachen} Crawl mit Tiefe 2 und 60 Seiten; das Ziel „möglichst vollständig“ erfordert dagegen Stunden. Die Kosten blieben trotzdem klein, weil Fargate nur während des Laufs abgerechnet wird.

\begin{table}[htbp]
\centering\small
\caption{Rechnung im MVP-Vortrag gegenüber den realen Großläufen.}\label{tab:rechnung-realitaet}
\begin{tabularx}{\linewidth}{@{}X l l l@{}}
\toprule
 & \textbf{MVP-Rechnung (Juli)} & \textbf{Aug-Lauf} & \textbf{Deep Run} \\
\midrule
Laufzeit gesamt & $\sim$ 10\,min (gerechnet) & \zDauerAug{}\,h (gemessen) & \zDauerDeep{}\,h (gemessen) \\
Hochschulen & 424 & 411 & alle Hochschulen, Breitensuche \\
Seiten je Hochschule & 60 & 800 & 2.500 \\
Rechenkosten (Fargate, 4\,vCPU/16\,GB) & $\sim$ 1,30\,\$ (Lambda) & $\sim$ 1,5\,\$ & $\sim$ 3\,\$ \\
Heruntergeladene PDFs & -- & 65.944 & 132.193 \\
\bottomrule
\end{tabularx}
\end{table}

\section{Infrastruktur als Code: Terraform, CLI, Bootstrap}\label{sec:aws-iac}
\paragraph{Terraform (11.\,Juli).} Zunächst wurde ein Terraform-Stack angelegt (EC2, Elastic IP, Security Group, IAM, S3, DynamoDB, ECR, Lambda). In der Praxis scheiterte der Weg an der Umgebung: Terraform und Docker waren in der Entwicklungsumgebung nicht ohne Weiteres nutzbar, der Terraform-Zustand und die handangelegten Ressourcen liefen auseinander, und die Wirkung des Stacks war schwer nachzuvollziehen. Der Stack wurde am 24.\,Juli als \emph{deprecated} markiert; CI/CD und alle späteren Ressourcen (Fargate, IAM) liefen über AWS-CLI und eingecheckte JSON-Vorlagen mit Platzhaltern (\texttt{ACCOUNT\_ID}, \texttt{REGION}).

\paragraph{Bootstrap-Skript (12.\,August).} Die Sandbox-Löschung (Abschnitt~\ref{sec:aws-verluste}) machte die Reproduzierbarkeit zum Hauptthema. Das Skript \texttt{webscraper/deploy/aws\_bootstrap.ps1} baut die gesamte Umgebung idempotent aus einer leeren Konto-ID auf: S3, DynamoDB, ECR, Lambda-Rolle, Image-Build und -Push, Lambda-Funktion, EC2 (IAM, Security Group, Key Pair, Instanz, Elastic IP; die Instanz konfiguriert sich über User-Data), GitHub-OIDC-Rolle, Fargate (Cluster, Rollen, Egress-Security-Group, Task-Definition). Am Ende gibt es die beiden GitHub-Variablen aus, die gesetzt werden müssen. Einige Lektionen sind im Skript fest verankert:
\begin{itemize}
  \item Die Richtlinien-Vorlagen müssen \emph{reine} Platzhalter enthalten. Zwei Vorlagen enthielten noch die alte Konto-ID (ein Leak im öffentlichen Repository, der zwar keinen Zugriff gewährt, aber behoben wurde).
  \item Ersetzungen laufen case-sensitiv (\texttt{-creplace}), sonst wird der kleingeschriebene Schlüssel \texttt{awslogs-region} der Task-Definition beschädigt.
  \item Das Skript muss reines ASCII sein, weil Windows PowerShell~5.1 Dateien ohne BOM als CP1252 liest und ein Gedankenstrich die Tokenisierung zerstört.
  \item Neue IAM-Rollen sind nicht sofort verwendbar; die Lambda-Erstellung wiederholt bis zu sechsmal im Abstand von zehn Sekunden.
\end{itemize}

\section{CI/CD mit GitHub Actions}
Drei Workflows liegen im Repository: \texttt{deploy-lambda.yml} (Image bauen, nach ECR pushen, Lambda aktualisieren, Pfadfilter auf \texttt{webscraper/} und \texttt{mlclassifier/}), \texttt{deploy-webapp.yml} (Instanz per Tag finden, per SSM-Run-Command \texttt{remote\_update.sh} ausführen: \texttt{git reset --hard}, \texttt{docker compose up -d --build}) und \texttt{bulk-crawl.yml} (Fargate-Task starten). Die Authentifizierung läuft per OIDC (kein statischer Schlüssel im Repository), die Anbindung per SSM (kein offener SSH-Port).

Beim ersten Einsatz am 22.\,Juli traten mehrere Fehler in kurzer Folge auf, jeweils mit einem eigenen Commit („fixed deploy“, „set home variable for git in ssm deploy“): Das Skript fand das App-Verzeichnis auf der Instanz nicht (Auto-Erkennung für drei mögliche Pfade nötig), und Git verlangte unter SSM eine gesetzte \texttt{HOME}-Variable. Für den interaktiven Zugang scheiterte \texttt{aws ssm start-session} zunächst, weil das Session-Manager-Plugin lokal fehlte.

\section{Betrieb der Webapp}
Die Webapp läuft als Docker-Compose-Stack (Webapp und Caddy) auf einer t3.small mit Elastic IP; die Datenbank liegt auf einem benannten Volume und überlebt Deployments. Daraus folgt, dass das Admin-Passwort das bei der Erstanlage gesetzte bleibt; ein nicht-destruktives Zurücksetzen per SQL-Update im Container wurde dokumentiert. Zwei Betriebsvorfälle sind erwähnenswert: Am 2.\,September war die Instanz per SSM nicht erreichbar („TargetNotConnected“, nach Ablauf der Sitzungsdaten), ein Reboot behob das. Außerdem erreichten neu eingeführte Umgebungsvariablen (Discovery, Rendering) den Container zunächst nicht, weil sie in der \texttt{docker-compose.yml} nicht durchgereicht wurden.

\section{Zugriffsbeschränkungen und ihre Folgen}\label{sec:aws-restriktionen}
\begin{itemize}
  \item \textbf{Bedrock-Mantle gesperrt.} Eine SCP verbietet \texttt{bedrock-mantle} (403), erlaubt aber das native \texttt{bedrock:InvokeModel}. Die LLM-Anbindung wurde daher auf den nativen Client umgestellt, mit dem EU-Inference-Profil \texttt{eu.anthropic.claude-haiku-4-5-20251001-v1:0}.
  \item \textbf{Keine langlebigen Schlüssel.} Der Account erlaubt nur kurzlebige Bedrock-API-Schlüssel (rund 12\,Stunden), die SSO-Sitzung verfällt teils alle 30 Minuten. Die LLM- und OCR-Läufe mussten deshalb in kleine, wiederaufnehmbare Blöcke („Chunks“) zerlegt werden; ein Skript (\texttt{next\_chunk.py}) schneidet nur den noch offenen Rest und überspringt bereits geprüfte URLs.
  \item \textbf{Kostensicht gesperrt.} Cost Explorer (\texttt{ce:GetCostAndUsage}) ist per SCP verboten; Kosten sind deshalb nur abgeschätzt (Abschnitt~\ref{sec:aws-kosten}).
  \item \textbf{TLS-Interception.} Im Firmennetz ersetzt ein Proxy Zertifikate; AWS-Aufrufe aus Python scheiterten mit \texttt{CERTIFICATE\_VERIFY\_FAILED}. Lokale Läufe nutzen \texttt{--no-verify-ssl}, die EC2-Instanz und Fargate dürfen die Prüfung nicht abschalten.
  \item \textbf{Bedrock-Kontingente.} Neue Konten melden zunächst 403 („account currently being verified“, nach unter zwei Stunden behoben). Am 30.\,September führte das Tageskontingent (\texttt{429 Too many tokens per day}) zum Abbruch der LLM-Läufe, die am Folgetag fortgesetzt wurden.
\end{itemize}

\section{Rückschläge und Verluste}\label{sec:aws-verluste}
\Cref{tab:vorfaelle} listet die schwerwiegenden Vorfälle. Die ersten beiden haben das Projekt am stärksten beeinflusst.

\begin{table}[htbp]
\centering\small
\caption{Betriebsvorfälle und Gegenmaßnahmen.}\label{tab:vorfaelle}
\begin{tabularx}{\linewidth}{@{}p{2.0cm} p{3.6cm} X@{}}
\toprule
\textbf{Datum} & \textbf{Ereignis} & \textbf{Folge und Gegenmaßnahme} \\
\midrule
12.08. & \textbf{Sandbox-Lease abgelaufen}, alle Ressourcen gelöscht, neuer Pool (Konto-ID geändert) & S3 (Manifeste, Feedback-Verdikte, Läufe), DynamoDB, Lambda, ECR, EC2 samt Webapp-Datenbank weg. Erhalten blieben Git-Repository, lokale Trainingsdaten und das Modell. Neuaufbau über das Bootstrap-Skript; der Großlauf über die Hochschulliste wurde noch am selben Tag in der neuen Umgebung wiederholt. \\
{[}früher{]} & Weiterer Lease-Verlust bzw.\ Kontowechsel \TODO{Datum und Umfang des ersten Verlusts ergänzen} & Hartkodierte Konto-IDs in Richtlinien; Vorlagen auf Platzhalter umgestellt. \\
bis 16.09. & \textbf{Budgetsperre} (50~USD) des Sandbox-Pools & Analyse: kein hängender Task, nur Dauerressource EC2 (t3.small) plus Läufe, Speicher und Bedrock; keine verwaisten EBS-Volumes, IPs oder NAT-Gateways. Gegenmaßnahme: dreistufige Terminierungsgarantie (siehe unten). \\
22.07. & CI/CD scheitert beim ersten Deployment (SSM, \texttt{HOME}, App-Pfad) & vier Korrektur-Commits innerhalb von rund 20 Minuten. \\
03.09. & DynamoDB-Visited-Store überspringt bereits gecrawlte Hochschulen & Discovery-Wiederholung fand nichts Neues; Seeds-only-Modus umgeht den Store gezielt. \\
17.09. & Zusammenfassung des Deep Runs (\texttt{summaries/}) nicht in S3 & Task-Rolle ohne Schreibrecht auf das Präfix; Per-Hochschul-Diagnose des Laufs ging verloren (aus den Manifesten rekonstruierbar). \\
16.09. & Evaluation überschreibt LLM-Urteile (Fehler im Code) & Alle bisherigen LLM-Ergebnisse wurden stillschweigend ignoriert; Korrektur in \texttt{rescore\_decisions}; ehrliche Abdeckung sank von 314 auf 303. \\
\bottomrule
\end{tabularx}
\end{table}

\paragraph{Der Sandbox-Verlust im Detail.} Am Vormittag des 12.\,August war die Lease abgelaufen, während die Evaluation des ersten Großlaufs über die 411 Hochschulen vorbereitet wurde. Verloren gingen nicht nur Rechenressourcen, sondern auch \emph{Daten}: der S3-Bucket mit allen Manifesten und den von Menschen vergebenen Verdikten der Feedback-Schleife sowie die Webapp-Datenbank. Dass Training und Code lokal beziehungsweise in Git lagen, begrenzte den Schaden. Die Konsequenz war doppelt: technisch die Idempotenz des Bootstrap-Skripts, organisatorisch die Regel, dass \emph{nichts Wertvolles nur in der Sandbox} liegt. Alle späteren Manifeste und Zwischenergebnisse (\texttt{*.jsonl}) liegen deshalb zusätzlich lokal. Der Großlauf, auf den sich die Evaluation bezog, lag in der verlorenen Umgebung und wurde noch am 12.\,August (13:17--19:35~UTC) in der neuen Sandbox wiederholt; die gesamte weitere Auswertung beruht auf diesem zweiten Lauf \TODO{bestätigen: erster Lauf vor dem 5.08. im alten Konto, Wiederholung lieferte ähnliche Zahlen (rund 12.000 Positive)}.

\paragraph{Was beim ersten Durchlauf des Bootstrap-Skripts schiefging.} Der Neuaufbau am 12.\,August lief nicht auf Anhieb durch, und jede Panne ist im Skript heute abgefangen:
\begin{itemize}
  \item \emph{PowerShell-Fehlerbehandlung:} Mit \texttt{\$ErrorActionPreference = "Stop"} bricht Windows PowerShell~5.1 bei \emph{jeder} stderr-Ausgabe eines nativen Programms ab, auch bei Exit-Code~0. Phase~8 stoppte deshalb ohne echten Fehler; die Fehlerprüfung läuft seitdem explizit über Exit-Codes.
  \item \emph{Elastic IP nicht verknüpft:} Das Zuweisen der Elastic IP schlug fehl, weil die Instanz noch nicht bereit war; wegen des geänderten Fehlermodus blieb das unbemerkt. Die Webapp war unter der Elastic IP nicht erreichbar, bis die Zuordnung von Hand nachgeholt wurde. Das Skript wartet nun auf \texttt{instance-running} und prüft die Zuordnung mit Wiederholungen.
  \item \emph{Fargate-Cluster fehlte:} Beim ersten Bulk-Start meldete die Webapp \texttt{ClusterNotFoundException}, weil die Fargate-Phase nicht vollständig durchgelaufen war. Ein Teil-Durchlauf (\texttt{-SkipImage -SkipEc2 -SkipCicd}) holte sie nach, was die Idempotenz bestätigt.
\end{itemize}

\paragraph{Die Budgetsperre im Detail.} Nach rund zwei Wochen Abwesenheit war der Pool eingefroren. Eine Prüfung aller Regionen ergab keinen aktiven Verursacher; das Budget war über viele Wochen aus der immer laufenden EC2-Instanz, den Fargate-Läufen, dem S3-Speicher und den ersten Bedrock-Läufen aufgelaufen. Als Konsequenz wurden die Bedrock-Aufrufe für die bereits vorhandenen Treffer ausgesetzt, und es wurde eine Terminierungsgarantie eingebaut, die ein „Hängen“ als Kostenursache ausschließt (Abschnitt~\ref{sec:aws-kosten}).

\section{Kosten und Kostenkontrolle}\label{sec:aws-kosten}
Da Cost Explorer gesperrt ist, sind die Kosten in \cref{fig:B02_kosten} \emph{Schätzungen} aus Laufzeiten und Preislisten; bei der Abgabe sollen sie durch Abrechnungswerte ersetzt werden \TODO{Kosten bei Betreuer/Account-Verwaltung erfragen; Werte in bericht/daten/kosten.csv ersetzen}. Als Preisbasis dienen: Fargate 4\,vCPU/16\,GB rund 0,23\,\$/h, t3.small samt EBS rund 17\,\$/Monat, S3 rund 0,024\,\$/GB und Monat, Claude Haiku 4.5 rund 1\,\$ je Million Eingabe- und 5\,\$ je Million Ausgabetoken (bei etwa 2.300 Eingabetoken je Dokument). Die geschätzten Gesamtkosten liegen bei etwa \zKostenGesamt{}\,\$.

\abb[0.85]{htbp}{B02_kosten}{Geschätzte Kosten nach Posten. Schraffierte Balken sind Schätzungen (Abrechnung ausstehend).}

Bemerkenswert ist die Rangfolge: Nicht der Crawl dominiert die Kosten, sondern die \emph{Qualitätssicherung} durch das LLM (rund \zKostenBedrock{}\,\$) und der dauerhaft laufende Webapp-Server (rund \zKostenEc{}\,\$). Die eigentlichen Großläufe kosten zusammen nur wenige Dollar. Das bestätigt die Grundidee des Aufbaus, Rechenlast kurzlebig und Dauerbetrieb klein zu halten, und zeigt zugleich, dass eine Dauerressource im Budget des Sandbox-Pools relevanter ist als jeder Crawl.

\paragraph{Terminierungsgarantie.} Die Budgetsperre führte zu einem dreischichtigen Schutz gegen unbegrenzte Laufzeit: (1)~\texttt{DOWNLOAD\_TIMEOUT} von 60\,s je Anfrage, (2)~\texttt{CLOSESPIDER\_TIMEOUT} als Zeitlimit je Hochschule und (3)~ein Wachhund (\texttt{bulk\_run.\_arm\_runtime\_watchdog}), ein Daemon-Timer, der nach \texttt{BULK\_MAX\_RUNTIME\_SECONDS} den Prozess mit \texttt{os.\_exit(2)} beendet, auch wenn der Reactor hängt. Dadurch endet der ECS-Task und die Abrechnung. Die Funktion wurde getestet und feuert zuverlässig.

\section{Lehren aus dem AWS-Betrieb}
\begin{enumerate}
  \item \textbf{Reproduzierbarkeit schlägt Eleganz.} Das Bootstrap-Skript aus AWS-CLI-Befehlen war robuster als der „saubere“ Terraform-Stack, weil es mit einem unbekannten Zielkonto umgehen konnte.
  \item \textbf{Daten sind wertvoller als Infrastruktur.} Infrastruktur lässt sich neu bauen, Labels und Zwischenergebnisse nicht. Wertvolles gehört lokal gesichert.
  \item \textbf{Dauerressourcen zählen.} Im Budgetrahmen einer Sandbox ist eine immer laufende Kleininstanz gewichtiger als ein zwölfstündiger Großlauf.
  \item \textbf{Zugriffsbeschränkungen sind Entwurfsvorgaben.} SCP-Sperren (Bedrock-Mantle, Cost Explorer), kurzlebige Schlüssel und TLS-Proxy erzwangen wiederaufnehmbare Verarbeitung in kleinen Blöcken.
  \item \textbf{Gerechnet ist nicht gemessen.} Die Lambda-Rechnung des MVPs („zehn Minuten“) war für das Zielbild um etwa den Faktor 40 zu optimistisch; erst Messungen am Großlauf haben die richtige Laufzeitklasse gezeigt.
\end{enumerate}

\chapter{Der Crawler: Entwicklung und Recall-Arbeit}\label{ch:crawler}

Der Crawler ist das Herzstück des Systems und zugleich das, was sich am häufigsten geändert hat. Dieses Kapitel folgt dem Weg von einem einfachen Tiefen-Crawl zu einer gesteuerten, seeds-gestützten Breitensuche. Es nennt zu jedem Schritt das Problem, die Maßnahme und die Wirkung, und es beschreibt ausdrücklich, was nicht funktioniert hat.

\section{Ausgangspunkt: ungesteuerter Tiefen-Crawl}
Der erste Spider folgte allen Links derselben Domain bis zu einer festen Tiefe (zunächst 2) und einer Seitengrenze (60) und lud PDF- und Word-Dateien herunter. Die Hochschulliste enthält nur Startseiten, daher wurde ein begrenzter Tiefen-Crawl nötig (\texttt{CRAWL\_MAX\_DEPTH}, \texttt{CRAWL\_MAX\_PAGES}, \texttt{allowed\_domains}). Im Juli zeigte ein Test auf der TU Dortmund das Grundproblem: Von der Startseite aus wurden 934 Dokumente gefunden, 372 heruntergeladen und \emph{nur 4} waren echte Modulhandbücher, bei rund 156 Studiengängen. Die Ursache war doppelt: Der Crawl folgte Links in beliebiger Reihenfolge und verbrauchte das kleine Seitenbudget für Nachrichten und Impressum, und Modulhandbücher liegen auf \emph{Fakultäts-Subdomains} (etwa \texttt{cs.tu-dortmund.de}) in Tiefen, die das Budget nie erreichte. Das Durchprobieren von Tiefe 1--3 und 25--200 Seiten in Lambda führte zudem ins Zeitlimit (Kapitel~\ref{ch:aws}).

\section{Fokussierter Crawl (24. Juli)}
Die erste größere Verbesserung bestand aus drei Maßnahmen:
\begin{enumerate}
  \item \textbf{Best-First-Frontier.} Jeder Link wird nach Schlüsselwörtern in URL und Ankertext bewertet (Umlaute gefaltet); die Bewertung wird zur Scrapy-Request-Priorität, Dokumente erhalten einen Grundbonus. Das Wortlisten-Vokabular (\texttt{modulhandbuch}, \texttt{modulbeschreibung}, \texttt{studiengang} usw. positiv; \texttt{impressum}, \texttt{news} usw. negativ) liegt im Profil.
  \item \textbf{Sitemap-Seeding.} \texttt{start\_requests} lädt \texttt{/sitemap.xml} sowie die in der robots.txt genannten Sitemaps, behandelt Gzip und Index-Rekursion und lädt passende Dokumente direkt.
  \item \textbf{Subdomain-Erkennung.} Springt ein Link in eine neue Subdomain derselben Basisdomain, wird die Tiefe zurückgesetzt, die Priorität erhöht und die Sitemap der Subdomain geladen (begrenzt über \texttt{CRAWL\_MAX\_SUBDOMAIN\_SITEMAPS}).
\end{enumerate}
Am selben Tag wurde die Engine \textbf{modularisiert}: Linkbewertung, Zieldefinition und Extraktion liegen in austauschbaren \emph{Extraktionsprofilen} (Kapitel~\ref{ch:architektur}). Außerdem wurden \textbf{skriptgenerierte PDFs} ohne \texttt{.pdf} in der URL erkannt (RWTH: \texttt{show\_document.asp?id=…}, \texttt{download.php?id=…}): Die Antwort wird anhand des Content-Types als Dokument erkannt und der Dateiname mit einem Hash versehen, damit sich Skript-URLs in S3 nicht überschreiben.

\subsection{Rauchtest auf vier Hochschulen}
Ein Test auf vier Hochschulen (769 PDFs; 202 positiv, 41 Review, 526 negativ) zeigte, dass die Wirksamkeit stark von der Seitenstruktur abhängt (\cref{tab:smoke}). Die Mechanismen funktionierten, der Recall ist \emph{strukturabhängig}: nahezu vollständig bei einfachen Seiten, teilweise bei Riesenuniversitäten und null bei JavaScript-Startseiten. Diese Erkenntnis bestimmte die weitere Arbeit.

\begin{table}[htbp]
\centering\small
\caption{Rauchtest auf vier Hochschulen (24.\,Juli).}\label{tab:smoke}
\begin{tabularx}{\linewidth}{@{}l l X@{}}
\toprule
\textbf{Hochschule} & \textbf{Ergebnis} & \textbf{Beobachtung} \\
\midrule
FH Wedel & sehr gut & Informatik-Handbuch mit Score 0,96 gefunden \\
TU Dortmund & teilweise & 8 Fakultäts-Subdomains erkannt, Informatik-Handbuch gefunden, aber nur etwa 28 von 156 Studiengängen (Budget) \\
RWTH Aachen & teilweise & skriptgenerierte PDFs erkannt (Review-Band um 0,6), die Informatik-Fakultät aber nicht erreicht \\
FH Münster & \textbf{null} & JavaScript-Startseite ohne Sitemap: Scrapy sieht keine Links \\
\bottomrule
\end{tabularx}
\end{table}

\section{Erster Großlauf über die Hochschulliste (12. August)}
Der Großlauf über 411 Listeneinträge (404 Hochschulen nach Zusammenführung von Einträgen mit gleicher Domain) lief \zDauerAug{}\,h auf einem Fargate-Task mit 4\,vCPU und 16\,GB und lud 65.944 PDFs herunter (Mittel rund \zRateAug{} PDFs pro Stunde, Median 285~s je Hochschule). Die Evaluation (eigenes Auswertungspaket mit HTML-Bericht, Szenarien nach Typ, Träger, Größe und Bundesland) ergab:
\begin{itemize}
  \item 11.579 automatisch positive, 8.166 Review-Fälle und 46.199 negative Dokumente;
  \item Abdeckung nur 261 von 404 Hochschulen, \textbf{145 Hochschulen ohne ein einziges Handbuch}, damit lag die Zahl der Treffer bei knapp 30\,\% des Erwartungswerts von rund \zZielwert{};
  \item Die Lücke war konzentriert auf kleine und nicht-öffentliche Hochschulen: öffentlich 79\,\%, privat 35\,\%, kirchlich 49\,\% Abdeckung; nach Größe von 49\,\% (unter 1.000 Studierende) bis 92\,\% (15.000--30.000);
  \item NORDAKADEMIE, bei der erwartungsgemäß keine Handbücher online liegen, lieferte 9 „positive“ Treffer; ein Hinweis auf Fehlalarme, der sich später bestätigte (Kapitel~\ref{ch:klassifikation});
  \item Die FernUniversität Hagen, mit Abstand die größte Hochschule Deutschlands, lieferte nur 57 Positive und 145 Review-Fälle; ein Hinweis auf zu geringe Tiefe pro Hochschule (siehe Phase~2 unten).
\end{itemize}

Die Diagnose zerlegte das Problem in drei unabhängige Lücken: (A)~82 Hochschulen wurden gecrawlt, hatten aber null Handbücher, (B)~70 Hochschulen lieferten null Dokumente, und (C)~auch bei Hochschulen mit Treffern war der Recall pro Hochschule gering. Stichproben zeigten, dass die als „leer“ gemeldeten Hochschulen (Konstanz, FH Kiel) sehr wohl Handbücher an erratbaren URLs haben (\texttt{/uploads/tx\_studiengang/modulebook\_*.pdf}, \texttt{/fileadmin/…/modulhandbuch\_*.pdf}): Nicht die Abwesenheit, sondern \emph{Auffindbarkeit} war der Engpass.

\section{Schnelle Gewinne ohne Zusatzbudget}
Zwei Anpassungen kosteten nichts:
\begin{itemize}
  \item \textbf{Schwelle:} Das ausgelieferte Modell hatte eine obere Schwelle von 0,65. Ein Neubewerten der \emph{bereits heruntergeladenen} 65.944 Dokumente bei 0,50 ergab +2.333 Positive (11.579 $\to$ 13.912) und +10 Hochschulen (261 $\to$ 271), bei rund 88\,\% Precision auf dem Hold-out. Ein „What-if“-Schalter der Evaluation (\texttt{--upper}) rechnet Entscheidungen aus gespeicherten Scores neu, ohne Neu-Crawl. (Dass diese Absenkung später zur Positiv-Flut beitrug, ist Teil von Kapitel~\ref{ch:klassifikation}.)
  \item \textbf{Vokabular:} Dateinamenmuster, die vorher ohne Gewicht blieben (\texttt{modulebook}, \texttt{mhb}, \texttt{handbuch}, \texttt{spo}, \texttt{studienverlaufsplan}), wurden ins Profil aufgenommen, sodass Best-First sie früher erreicht.
\end{itemize}

\section{Discovery: Seeds aus Suchindizes}\label{sec:crawler-discovery}
Dominanter Hebel war die \emph{Discovery}: Pro Hochschule wird vorab gesucht, wo Modulhandbücher liegen, und diese URLs werden als priorisierte Ziele in den Crawl eingespeist (\texttt{DISCOVERY\_SEEDS\_PATH}, in der Cloud über S3). \Cref{tab:discovery} zeigt alle erprobten Anbieter, auch die verworfenen.

\begin{table}[htbp]
\centering\small
\caption{Erprobte Discovery-Quellen.}\label{tab:discovery}
\begin{tabularx}{\linewidth}{@{}l l X@{}}
\toprule
\textbf{Quelle} & \textbf{Status} & \textbf{Befund} \\
\midrule
Common Crawl (CDX) & \textbf{genutzt} & Kostenlos, ohne Schlüssel; je Domain PDFs mit „modul/handbuch“-Mustern. Von 18 getesteten Problemdomänen wurden für 11 Handbücher gefunden; stark bei Fakultäts- und Modul-Datenbank-Subdomains. Blind für tief verlinkte, unverlinkte Dateien. Mehrere Monatsindizes werden vereinigt, Anfragen mit Backoff und Pausen (CC drosselt Bursts und sperrte die Test-IP zeitweise). \\
Serper (Google-SERP) & \textbf{genutzt} & Abfrage \texttt{site:<basis> <begriff>}; findet tiefe Dateien wie das Konstanzer \texttt{modulebook}. Die kostenlose Stufe verlangt Anpassungen (kein \texttt{filetype:}, höchstens 10 Treffer je Anfrage); eine bezahlte Stufe war nicht möglich. \\
Firecrawl & vorbereitet, kaum getestet & Rendert JavaScript, liefert Link-Karten und ersetzt eine Seite durch ihr gerendertes HTML (Middleware). Für IU hat die Karte \emph{keine} PDF-Handbücher gefunden, das Portal ist ein reiner JavaScript-Katalog. Die Pipeline für Katalog-Seiten (Inhaltsklassifikation statt PDF) existiert, wurde aber nicht im Großlauf eingesetzt. \\
\textcolor{orange!80!black}{Google Custom Search} & \textbf{verworfen} & Google hat „gesamtes Web durchsuchen“ aus der Programmable Search Engine entfernt; nur noch 50 gelistete Domains, damit für 400 Hochschulen nutzlos. \\
\textcolor{orange!80!black}{DuckDuckGo} & \textbf{verworfen} & Aus der Cloud-Umgebung CAPTCHA-gesperrt. \\
\textcolor{orange!80!black}{Sitemaps allein} & unzureichend & Tiefe Dateien (z.\,B.\ Konstanz) stehen in keiner Sitemap; die Universität hat keine \texttt{sitemap.xml}. \\
Brave Search API & nicht gebaut & Unabhängiger Index, wäre die einzige realistische Quelle für die tiefsten Dateien, benötigt aber Anmeldung und eventuell Zahlungsmittel. \\
\bottomrule
\end{tabularx}
\end{table}

\paragraph{Wirkung.} Ein Test auf drei Hochschulen, die im ersten Lauf null Treffer hatten (htwsaar, Konstanz, FH Kiel), ergab mit Serper und Common Crawl nach einem kurzen Crawl 42 klassifizierte Dokumente, davon 7 automatisch positive; die LLM-Zweitstufe bestätigte weitere 6 englische Handbücher, die das TF-IDF-Modell unsicher beurteilt hatte. Zusammen 13 Handbücher bei vorher null.

\section{Seeds-only-Modus: Lücken ohne Duplikate schließen}
Der Visited-Store in DynamoDB dedupliziert auf Ebene der \emph{Hochschule} (Start-URL), nicht des Dokuments: Ein erneuter Lauf mit Discovery übersprang jede bereits gecrawlte Hochschule, und die Discovery lief nie. Ein vollständiges Leeren der Tabelle hätte alle Hochschulen neu gecrawlt (und zu physischen Duplikaten in S3 geführt, da der S3-Schlüssel die Job-ID enthält). Die Lösung war ein \textbf{Seeds-only-Modus}: Der Spider holt ausschließlich die Seeds, überspringt den Visited-Store, und Hochschulen ohne Seeds werden sofort übergangen. Der erste Lauf (3.\,September) holte 431 PDFs von 134 Hosts (226 positiv, 179 negativ, 26 Review) und hob die Abdeckung \emph{in wenigen Minuten} von 271 auf 314 Hochschulen (+43), die Positive von 13.912 auf 14.133. Das bewies den Ansatz.

\subsection{robots.txt und Seeds}\label{sec:crawler-robots}
Bei elf Hochschulen blieb der Seeds-Lauf zunächst bei null heruntergeladenen Dateien, obwohl Seeds vorlagen. Die Vermutung war JavaScript oder robots.txt; tatsächlich lag es an einer veralteten Seed-Datei im Lauf. Das echte robots-Problem trat bei der Universität Konstanz auf (\texttt{Disallow: /}). Die Lösung ist \texttt{SEEDS\_BYPASS\_ROBOTS} (standardmäßig \emph{aus}): Es setzt \texttt{dont\_obey\_robotstxt} \emph{nur} auf Requests für Discovery-Seeds, also auf einzelne, vom Suchindex als öffentlich ausgewiesene Handbuch-URLs, während der Link-Crawl der Seite robots.txt weiter respektiert. Im Test lieferte Konstanz damit 13 von 18 Dateien (6\,MB). Diese Ausnahme ist eine bewusste Abwägung, die im Ausblick und in der Diskussion (Kapitel~\ref{ch:diskussion}) aufgegriffen wird.

\section{Phase 2: Vollständigkeit je Hochschule}
Auch bei Hochschulen mit Treffern blieb der Recall niedrig. Beispiel FernUniversität Hagen: nur 70 Handbücher, erwartet waren hunderte. Die Ursache: Der gesteuerte Crawl \emph{tunnelte} in die erste starke Fakultät (\texttt{/mi/}) und erreichte nie Fakultäten, deren Modulkataloge hinter neutraler Navigation liegen (\texttt{/psychologie/studium/portale/<abschluss>/…}, kein „modul“ im Pfad). Gegenmaßnahmen:
\begin{itemize}
  \item \textbf{Breadth-Fairness:} \texttt{CRAWL\_MAX\_PAGES\_PER\_SECTION} begrenzt die Seitenfolgen pro erstem Pfadsegment (Fakultät) und erzwingt Streuung; Dokument-Downloads werden nie gekappt.
  \item \textbf{Navigations-Tokens} im Profil (\texttt{portal}, \texttt{bsc}, \texttt{msc}, \texttt{studienangebot} usw.), damit schlüsselwortlose Programm-Zweige verfolgt werden.
  \item \textbf{Deep-Run-Konfiguration:} Tiefe 5, 2.500 Seiten, 60 Seiten pro Sektion, bis 3.000 Seiten und 1.800\,s je Hochschule, Gesamtlaufzeit bis 12\,h, \texttt{force}.
  \item \textbf{Terminierungsgarantie} und Diagnose-Statistiken je Hochschule (Abbruchgrund, Seitenzahl, Sektionsverteilung).
\end{itemize}
Der lokale Test an Hagen zeigte, dass der Crawl jetzt Wirtschaft, Psychologie (0 $\to$ 4 Positive) und Kultur- und Sozialwissenschaften erreichte statt nur einer Fakultät.

\subsection{Deep Run (16./17. September)}
Der Deep Run lief \zDauerDeep{}\,h und lud 132.193 PDFs herunter (37.993 automatisch positiv für sich genommen). Hagen stieg von 70 auf \textbf{870} Positive (Faktor 12), verteilt über sieben Fakultäten. Die kumulierte Vereinigung aller Läufe (Aug-Lauf, Seeds-Läufe, Deep Run) ergab \zDocsTotal{} Dokumente, 43.270 TF-IDF-Positive bei Schwelle 0,5 und eine Abdeckung von 321 Hochschulen. Das erreichte \emph{nominell} das 40.000er-Ziel, war aber durch die Precision-Probleme noch nicht belastbar (Kapitel~\ref{ch:klassifikation}).

\abb[0.95]{htbp}{B03_verlauf_abdeckung}{Abdeckung (links) und gelistete Positive (rechts) im Projektverlauf. Die grauen Balken sind TF-IDF-Positive ohne LLM-Prüfung; erst die Endstände (grün, blau) sind qualitätsgesichert.}

\section{Höflichkeit und Grenzen}
Der Crawler respektiert robots.txt, drosselt (Basiswartezeit 1\,s plus AutoThrottle) und tritt mit identifizierbarem User-Agent auf. Weitere harte Grenzen (Dateigröße 50\,MB, Timeout 60\,s je Anfrage) sind in Anhang~\ref{app:konfig} aufgeführt. Weiterhin außerhalb der Reichweite sind Handbücher hinter Login, nur zur Laufzeit erzeugte Dokumente ohne Datei und Handbücher auf einer \emph{anderen} registrierbaren Domain als der Startseite (\texttt{\_same\_site} vergleicht die letzten zwei DNS-Labels).

\section{Zusammenfassung der Crawl-Entwicklung}
\Cref{tab:crawl-weg} bündelt den Weg.

\begin{table}[htbp]
\centering\small
\caption{Crawl-Strategien im Zeitverlauf mit Problem, Maßnahme und Wirkung.}\label{tab:crawl-weg}
\begin{tabularx}{\linewidth}{@{}p{2.2cm} p{3.6cm} X@{}}
\toprule
\textbf{Zeitpunkt} & \textbf{Problem} & \textbf{Maßnahme und Wirkung} \\
\midrule
Juli & Budget für Unwichtiges verbraucht, TU Dortmund 4 Treffer in 934 Dokumenten & Best-First, Sitemap, Subdomains; Handbuch-Funde deutlich höher \\
24.07. & Skript-URLs, JS-Startseiten & Content-Type-Erkennung (RWTH gelöst); JS-Startseiten (FH Münster) \emph{ungelöst} \\
12.08. & 145 Hochschulen ohne Treffer & Auswertungsmodul; Schwelle 0,5: +2.333 Positive, +10 Hochschulen \\
27.08. & Handbücher nicht auffindbar & Discovery (Common Crawl, Serper), OCR, LLM-Zweitstufe \\
03.09. & Visited-Store blockiert Wiederholung & Seeds-only: +43 Hochschulen in Minuten \\
16.09. & robots-Sperre (Konstanz), Tunneln (Hagen) & Seeds-Robots-Bypass; Breadth-Fairness; Deep Run: Hagen 70 $\to$ 870 \\
\bottomrule
\end{tabularx}
\end{table}

\chapter{Klassifikation: vom TF-IDF-Modell zur LLM-Zweitstufe}\label{ch:klassifikation}

Dieses Kapitel ist der Kern des Machine-Learning-Anteils. Es zeigt, wie ein auf dem Papier sehr gutes Modell im realen Crawl deutlich schlechter abschnitt, wie das erkannt wurde und welche Maßnahmen die Qualität schließlich auf ein belastbares Niveau gebracht haben.

\section{Aufgabe und Daten}
Die Aufgabe ist binär: \emph{Ist dieses PDF ein Modulhandbuch?} Als Modulhandbuch zählt (Definition im Prompt der Zweitstufe) ein systematisches Handbuch, das die Module eines Studiengangs beschreibt. Einzelne Modulblätter, Prüfungsordnungen, Studienverlaufspläne und Vorlesungsverzeichnisse zählen \emph{nicht}. Diese strenge Definition ist eine Entscheidung; sie ist der Hauptgrund für viele Beinahe-Treffer.

Die Trainingsdaten (\texttt{modulhandbuecher/\{positiv,negativ\}/<hochschule>/}) wuchsen in drei Stufen (\cref{tab:modelle}):
\begin{itemize}
  \item \textbf{MVP-Modell (Juli):} 202 Dokumente (134 positiv, 68 negativ) aus 21 Hochschulen nach Deduplizierung per normalisiertem Text-Hash vor dem Split.
  \item \textbf{Nach Feedback (22.\,Juli):} Die Review-UI lieferte 157 Verdikte (9 positiv, 148 negativ), das nachtrainierte Modell nutzte 356 Dokumente aus 22 Hochschulen (143 positiv, 213 negativ). Die Verdikte waren vor allem harte Negative, das heißt Dokumente, die dem Modell zuvor als Handbuch erschienen.
  \item \textbf{Nach der Massenerhebung (3.\,September):} Mit über 5.000 zusätzlichen Handbüchern aus dem Crawl wuchs der Datensatz auf \textbf{2.695 Dokumente} aus \textbf{269 Hochschulen} (1.971 positiv, 724 negativ).
\end{itemize}

\section{Modell und Evaluation}
Das Modell ist eine logistische Regression auf TF-IDF-Merkmalen (Wort- und Zeichen-n-Gramme) des extrahierten PDF-Texts. Die Evaluation nutzt \emph{StratifiedGroupKFold, gruppiert nach Hochschule}: Keine Hochschule steht gleichzeitig im Training und im Test, es wird also die Erkennung an \emph{ungesehenen} Hochschulen gemessen (Szenario~B). Ein Leakage-Check trainiert zwei Varianten, \emph{nur Inhalt} und \emph{Inhalt plus Dateiname/Metadaten}; sind sie ähnlich gut, liest das Modell das Dokument und nicht den Dateinamen. Die Schwellen werden auf 95\,\% Recall kalibriert (Kapitel~\ref{ch:grundlagen}).

\begin{table}[htbp]
\centering\small
\caption{Modellversionen und gruppierte Kreuzvalidierung.}\label{tab:modelle}
\begin{tabular}{@{}l r r r r r r@{}}
\toprule
\textbf{Version} & \textbf{Dok.} & \textbf{Hochsch.} & \textbf{PR-AUC} & \textbf{ROC-AUC} & \textbf{Prec.} & \textbf{Recall} \\
\midrule
MVP (Juli), Variante \emph{content} & 202 & 21 & 0,957 & 0,941 & 0,871 & 0,955 \\
MVP, \emph{content + Dateiname} & 202 & 21 & 0,960 & 0,926 & 0,848 & 0,955 \\
Stand September, \emph{content} & 2.695 & 269 & 0,988 & 0,973 & 0,965 & 0,950 \\
September, \emph{content + Metadaten} & 2.695 & 269 & 0,976 & 0,959 & 0,958 & 0,950 \\
\bottomrule
\end{tabular}
\par\smallskip\footnotesize Precision und Recall am Schwellenwert für das Recall-Ziel von 95\,\% (MVP: 0,387; September: 0,627). Beim Schwellenwert 0,5 erreicht die September-Variante \emph{content} Precision 0,952 und Recall 0,985.
\end{table}

Die Leakage-Prüfung bestätigte das Modell: Inhalt und Inhalt-plus-Metadaten liegen nahe beieinander (MVP 0,957 gegenüber 0,960; September sogar leicht schlechter mit Metadaten), also wird das schlankere Inhaltsmodell ausgeliefert. Die Konfusionsmatrix der September-Variante am Recall-Ziel: 1.873 richtig positive, 98 falsch negative, 656 richtig negative, 68 falsch positive. Anhang~\ref{app:metriken} enthält die Details.

\section{Der Praxis-Gap: 96,5\,\% im Training, 25\,\% im Crawl}
Die Trainingsmetriken waren hervorragend, doch im realen Crawl verhielt sich das Modell anders. Zwei Beobachtungen kamen aus dem Betrieb:
\begin{enumerate}
  \item NORDAKADEMIE lieferte neun „Handbücher“, obwohl dort keine online sind.
  \item Die LLM-Verifikation der 226 automatisch positiven Treffer aus dem Seeds-Lauf (215 prüfbar) ergab nur \textbf{129 echte Handbücher} und 86 Fehlalarme (rund 40\,\%): Studienordnungen, Verkündungsblätter, ein Bewerbungsformular mit Score 0,79, Lehrbücher.
\end{enumerate}
Die systematische Auswertung nach Score-Bereichen (\cref{fig:03_precision_je_score}) zeigt, wie stark der Score überschätzt: Bei Score 0,5--0,6 bestätigt das LLM nur 8\,\% der Treffer, bei 0,8--0,9 sind es 44\,\%; erst oberhalb von 0,9 liegt die Bestätigungsrate bei \zPilotPct{}\,\% (Pilotstichprobe, $n=\zPilotN{}$, Wilson-Intervall \zPilotLo{}--\zPilotHi{}\,\%).

\abb[0.8]{htbp}{03_precision_je_score}{Anteil der vom LLM bestätigten TF-IDF-Positive je Score-Bereich (LLM als Referenz). Die Bereiche oberhalb 0,9 stammen aus der Pilotstichprobe.}
\abb[0.8]{htbp}{23_training_vs_praxis}{Precision im Training (gruppierte Kreuzvalidierung) gegenüber dem realen Crawl. Die Differenz ist die Kernaussage zur Übertragbarkeit.}

\paragraph{Ursachen.} (1)~Das Trainingsset bildet die Breite der Beinahe-Treffer im Web nicht ab: Die Negativbeispiele waren überwiegend leichte Fälle, während der Crawl Zehntausende Prüfungsordnungen, Modulblätter und Studienverlaufspläne enthält. (2)~Die Schwelle wurde im Deep Run auf 0,5 gesenkt, um Recall zu gewinnen, das aktuelle Modell schlägt 0,8 vor; das erzeugte die Positiv-Flut (37.993 Positive allein im Deep Run). (3)~Die Verteilung unterscheidet sich zwischen Hochschulen mit Handbuch-Datenbank (viele ähnliche Dokumente je Host) und den Trainingshochschulen.

\paragraph{Welche Dokumente werden verwechselt?} Die LLM-Begründungen der \zFPTotal{} verworfenen TF-IDF-Positive wurden per Stichwort in Fehlerarten eingeteilt (\cref{fig:07_fehlerarten}): knapp die Hälfte sind Prüfungs- oder Studienordnungen (48\,\%), 19\,\% einzelne Modulblätter oder Syllabi, 12\,\% Studienverlaufspläne, 5\,\% unvollständige Modulkataloge, 4\,\% Veranstaltungsverzeichnisse und 2\,\% Formulare. Das sind genau die Dokumenttypen, die dem Modulhandbuch semantisch am nächsten liegen. Die Zuordnung ist stichwortbasiert und daher grob.

\abb[0.85]{htbp}{07_fehlerarten}{Fehlerarten der TF-IDF-Positive, die das LLM verworfen hat.}

\paragraph{Dateiname als Signal.} Obwohl das Modell nicht auf Dateinamen trainiert wurde, ist der Dateiname ein starkes Zusatzsignal (\cref{fig:08_dateiname_signal}): Enthält er „modul“, „handbuch“ oder „mhb“, werden bei Score $\ge$ 0,9 90\,\% bestätigt, ohne diese Stichwörter nur 41\,\%; im Bereich 0,5--0,7 sind es 36\,\% gegenüber 8\,\%. Das ist ein Ansatzpunkt für ein künftiges Modell (Kapitel~\ref{ch:fazit}).

\abb[0.85]{htbp}{08_dateiname_signal}{Bestätigungsrate des LLM mit und ohne Handbuch-Stichwort im Dateinamen.}

\section{Human-in-the-Loop}
Die erste Antwort auf Modellfehler war ein \emph{Menschen-Review}. Die Review-UI der Webapp zeigt relevante Dokumente (Positive und Review-Band) nach Score sortiert, mit Seitenweise-Anzeige für große Mengen und Download. Verdikte werden gesammelt (nicht einzeln) an S3 geschrieben, über das Kommando \texttt{feedback-retrain --from-s3} geholt, die Dokumente unter dem \emph{menschlichen} Label in den Trainingsbaum einsortiert (nach Hochschule gruppiert) und das Modell neu trainiert; das neue Modell wird per Git eingecheckt und über CI/CD ausgeliefert. Der Mechanismus ist modellunabhängig (positiv/negativ), nicht modulhandbuchspezifisch. In der Praxis stieß der manuelle Weg an die Skala: Bei zehntausenden Kandidaten ist menschliches Prüfen nicht tragfähig; das war der Anlass für die LLM-Zweitstufe.

\section{LLM-Zweitstufe}\label{sec:llm}
\paragraph{Aufbau.} Die Zweitstufe (\texttt{mlclassifier/llm\_review.py}) liest die Manifest-Einträge eines Bandes (Review, automatisch positiv oder ein Score-Fenster), lädt die PDFs aus S3, extrahiert Text, sendet Textauszüge an Claude Haiku~4.5 auf Amazon Bedrock und erhält ein strukturiertes Urteil mit Begründung und Konfidenz zurück. Der Prompt wird aus einer \texttt{ClassificationSpec} gebaut (Definition, Beispiele, Ausschlüsse), die Zweitstufe ist also wiederverwendbar. Ein Urteil gilt nur bei Konfidenz $\ge$ 0,8 als bestätigt (praktisch liegen alle Urteile zwischen 0,8 und 1,0). Die Verarbeitung läuft parallel (Thread-Pool, standardmäßig 12 Arbeiter, mit Sperre für Manifestschreiben und fehlerisolierten Einzelaufrufen); das senkte die Dauer für rund 6.000--8.000 Dokumente von etwa fünf Stunden auf etwa 30 Minuten. Das Modell wurde bewusst klein gewählt (Haiku): billig genug für zehntausende Dokumente, im Test aber hinreichend genau (6 von 6 Zweifelsfällen aus htwsaar mit Konfidenz 0,95 korrekt als Handbuch erkannt).

\paragraph{Läufe.} \Cref{tab:llm-laeufe} listet die Läufe.

\begin{table}[htbp]
\centering\small
\caption{LLM-Läufe der Zweitstufe.}\label{tab:llm-laeufe}
\begin{tabular}{@{}l r r r@{}}
\toprule
\textbf{Lauf} & \textbf{geprüft} & \textbf{bestätigt} & \textbf{Rate} \\
\midrule
Review-Band des Aug-Laufs (Score 0,3--0,7) & 6.065 & 1.845 & 30\,\% \\
Positive 0,5--0,9 des Deep Runs & 20.966 & 5.153 & 25\,\% \\
Pilot: Stichprobe der Positive $\ge$ 0,9 & 400 & 330 & 82,5\,\% \\
OCR-Treffer (Score $\ge$ 0,5) & 247 & 42 & 17\,\% \\
\midrule
\textbf{Summe (je URL eindeutig)} & \textbf{\zLLMReviewed{}} & \textbf{\zLLMConfirmed{}} & \\
\bottomrule
\end{tabular}
\end{table}

\paragraph{Die Pilotstichprobe.} Um nicht auch die 13.741 Dokumente oberhalb von 0,9 mit dem LLM prüfen zu müssen (rund 35\,\$), wurde eine Zufallsstichprobe von 400 geprüft. 330 wurden bestätigt (\zPilotPct{}\,\%, Wilson-Intervall \zPilotLo{}--\zPilotHi{}\,\%). Daraus folgt eine \emph{Schätzung} von rund \zEstTrue{} echten Modulhandbüchern in der erweiterten Liste (Bereich \zEstTrueLow{}--\zEstTrueHigh{}). Die Annahme ist, dass die Stichprobe für die gesamte Gruppe repräsentativ ist; das ist nicht für alle Hosts erfüllt (siehe Abschnitt~\ref{sec:grenzen}).

\abb[0.9]{htbp}{05_llm_laeufe}{LLM-Läufe: Umfang und Bestätigungsrate.}

\paragraph{Ein Fehler in der Auswertung.} Ein Fehler in der Evaluation (\texttt{rescore\_decisions}) überschrieb jedes LLM-Urteil beim Neubewerten mit der Score-Schwelle: Ein vom LLM abgelehnter Fehlalarm kippte zurück auf „positiv“, sobald sein Score $\ge$ 0,5 war. Damit waren \emph{alle} bisherigen LLM-Ergebnisse in der Auswertung stillschweigend wirkungslos. Der Fix: Datensätze mit \texttt{llm\_reviewed=True} werden nie neu bewertet, das LLM-Urteil ist Wahrheit. Wirkung: positive Dokumente 14.133 $\to$ 13.565, Review 5.852 $\to$ 2.122, Abdeckung 314 $\to$ 303. Die \emph{ehrliche} Zahl war also kleiner als die zuvor berichtete.

\section{OCR für gescannte PDFs}
Rund 3.600 PDFs lieferten keinen Text (\texttt{empty\_document}), landeten im Review-Band und sollten nie stillschweigend im Negativ verschwinden. Ein OCR-Fallback (PyMuPDF rastert, Tesseract mit \texttt{deu+eng}) wurde gebaut und lokal auf \zOCRDone{} Scans angewendet: \zOCRText{} lieferten Text, nur \zOCRLLM{} erreichten einen Score $\ge$ 0,5 und gingen ans LLM, \zOCRConf{} wurden bestätigt (\cref{fig:09_ocr_trichter}). Der Ertrag (rund 1\,\% der Ausgangsmenge) steht in keinem Verhältnis zum Aufwand (lokale Tesseract-Installation, mehrere parallel laufende Prozesse über viele Stunden); ein klassisches Beispiel für einen Ansatz, der gebaut wurde, weil er plausibel war, und dessen Grenznutzen erst die Messung zeigte.

\abb[0.8]{htbp}{09_ocr_trichter}{OCR-Trichter: von den Scans ohne Textebene bis zu den bestätigten Handbüchern.}

\section{Erwogene Alternativen}
\begin{itemize}
  \item \textbf{Regex auf „Modulhandbuch“ im Dateinamen/Text:} scheitert in beide Richtungen (Handbücher ohne das Stichwort im Namen, Verordnungen und Formulare, die „Modulhandbuch“ erwähnen).
  \item \textbf{Transformer-Finetuning (z.\,B.\ BERT):} bei 202 bzw.\ 356 Dokumenten ungeeignet; bei 2.695 Dokumenten denkbar, aber die LLM-Zweitstufe liefert dasselbe ohne Training.
  \item \textbf{LLM für alle Dokumente:} bei 164.396 PDFs und rund 2.300 Eingabetoken je Dokument wären das gut 375 Millionen Token; mit dem TF-IDF-Filter genügen rund 27.700 Dokumente. Die zweistufige Architektur spart gut 80\,\% der LLM-Aufrufe.
  \item \textbf{Rein manuelles Prüfen:} Bei 30\,s je Dokument wären das rund \zManuellOhneH{}\,Stunden für die LLM-geprüften Dokumente allein (\cref{fig:B09_manueller_aufwand}).
\end{itemize}

\chapter{Ergebnisse und Kennzahlen}\label{ch:ergebnisse}

Dieses Kapitel beantwortet die Leitfragen aus Kapitel~\ref{ch:einleitung} mit Zahlen. Alle Werte stammen aus der Abschlussauswertung (\texttt{auswertung/data/kennzahlen.json}) und den S3-Auflistungen und sind im Bericht als Makros hinterlegt, sodass sie nach den abschließenden LLM-Prüfungen automatisch aktualisiert werden. \TODO{Nach Abschluss der laufenden LLM-Prüfung der ca.\ 6.000 verdächtigen Dokumente: build.ps1 ausführen, Zahlen und Text prüfen}.

\section{Kennzahlen auf einen Blick}
\begin{table}[htbp]
\centering\small
\caption{Zentrale Kennzahlen des Projekts.}\label{tab:kennzahlen}
\begin{tabularx}{\linewidth}{@{}X r@{}}
\toprule
\textbf{Kennzahl} & \textbf{Wert} \\
\midrule
Hochschulen in der Hochschulliste & \zUnisTotal \\
Heruntergeladene PDFs (URL-eindeutig, alle Läufe) & \zDocsTotal \\
Objekte / Datenmenge in S3 (inkl.\ physischer Duplikate) & \zSpeicherObjekte{} / \zSpeicherGB{}\,GB \\
Laufzeit Großlauf August / Deep Run & \zDauerAug{}\,h / \zDauerDeep{}\,h \\
Durchsatz August / Deep Run (PDFs pro Stunde im Mittel) & \zRateAug{} / \zRateDeep \\
TF-IDF-Positive vor jeder Prüfung (Schwelle 0,5) & \zInitialPositive \\
Vom LLM geprüfte Dokumente & \zLLMReviewed \\
\quad davon bestätigt / verworfen & \zLLMConfirmed{} / \zLLMRejected \\
Endliste streng (nur LLM-bestätigt) & \zFinalStrict \\
Endliste erweitert (zusätzlich TF-IDF $\ge$ 0,9, ungeprüft) & \zFinalListed \\
Geschätzt echte Modulhandbücher in der erweiterten Liste & \zEstTrue{} (\zEstTrueLow--\zEstTrueHigh) \\
Eindeutig nach Host und Dateiname (streng / erweitert) & \zUniqueStrict{} / \zUniqueLikely \\
Hochschulen mit $\ge$ 1 Handbuch (streng / erweitert) & \zUnisStrict{} (\zUnisStrictPct\,\%) / \zUnisLikely{} (\zUnisLikelyPct\,\%) \\
Hochschulen ohne jedes Dokument bzw.\ ohne Fund & \zUnisNone \\
Median Handbücher je Hochschule (erweitert) & \zMedianMH \\
Anteil der zehn ergiebigsten Hochschulen am Ergebnis & \zTopTenShare\,\% \\
\bottomrule
\end{tabularx}
\end{table}

\abb[0.9]{htbp}{01_trichter}{Trichter vom Crawl (\zDocsTotal{} PDFs) über die TF-IDF-Positive und die LLM-Prüfung bis zu den geschätzt echten Handbüchern.}

\section{Wie lange dauert ein Scraping-Lauf?}
\Cref{fig:B04_laeufe_dauer} und \cref{tab:rechnung-realitaet} zeigen die Läufe. Der Großlauf im August dauerte \zDauerAug{}\,h für 65.944 PDFs (rund \zRateAug{} PDFs pro Stunde, im Spitzenbereich bis rund 11.900 pro Stunde), der Deep Run \zDauerDeep{}\,h für 132.193 PDFs (rund \zRateDeep{} PDFs pro Stunde, Spitze rund 14.400). Der Seeds-only-Lauf holte 431 PDFs von 134 Hosts in unter vier Minuten. Der Median pro Hochschule lag im August bei 285 Sekunden. Entscheidend ist, dass die Läufe von wenigen großen Hochschulen dominiert werden, die Mehrzahl der Hochschulen war in Minuten abgearbeitet.

\abb[0.95]{htbp}{B04_laeufe_dauer}{Laufzeit, Dokumentenzahl und Durchsatz der drei Läufe.}
\abb[0.85]{htbp}{18_durchsatz}{Durchsatz (neue PDFs pro Stunde) im Verlauf der beiden Großläufe.}

\section{Wie viele Dokumente werden gefunden, und was geschieht mit ihnen?}
\Cref{fig:02_endstatus_nach_quelle} zeigt den Endstatus aller \zDocsTotal{} Dokumente nach Lauf: LLM bestätigt (\zLLMConfirmed), TF-IDF $\ge$ 0,9 ungeprüft (\zTfidfHigh), Positive aus anderen Läufen ungeprüft (\zTfidfOther), LLM verworfen (\zLLMRejected), TF-IDF negativ ungeprüft (\zTfidfNegative) und ungeklärt wegen fehlendem Text oder Fehlern (\zUnresolved). Die beiden Großläufe überlappen: 33.883 URLs kamen in beiden vor; im Deep Run gilt dann der jüngere Datensatz. Der Deep Run erreichte 718 Hosts, die der August-Lauf nicht kannte, der August-Lauf 296 Hosts, die der Deep Run nicht erreichte (\cref{fig:20_laeufe_vergleich}).

\abb[0.9]{htbp}{02_endstatus_nach_quelle}{Endstatus aller Dokumente nach Lauf.}
\abb[0.85]{htbp}{20_laeufe_vergleich}{Vergleich von August-Lauf und Deep Run (URL-Überlappung).}

\section{Wie groß sind die gespeicherten Datenmengen?}
Der S3-Bucket enthält zum Auswertungszeitpunkt \zSpeicherObjekte{} Objekte mit zusammen \zSpeicherGB{}\,GB (\zSpeicherGiB{}\,GiB); davon entfallen rund 55\,GB auf den August-Lauf und rund 99\,GB auf den Deep Run, die Seeds-Läufe wiegen unter 1\,GB (\cref{fig:B05_datenmenge}). Da erzwungene Wiederholungsläufe dieselben PDFs unter neuer Job-ID ablegen, sind rund \zDupVolumePct\,\% des Volumens physische Duplikate; die URL-eindeutigen Dokumente belegen \zUniqueGB{}\,GB. Ein durchschnittliches PDF ist \zMeanMB{}\,MB groß, der Median liegt bei \zMedianMBAll{}\,MB; echte Modulhandbücher sind mit einem Median von \zMedianMBMH{}\,MB größer (typischerweise 100~KB bis 10~MB). Die gelisteten Modulhandbücher (erweiterte Liste) belegen zusammen rund \zMHVolumeGB{}\,GB, die streng bestätigten rund \zMHStrictVolumeGB{}\,GB. Die Manifeste (Metadaten, Scores) wiegen insgesamt nur rund 0,2\,GB.

\abb[0.95]{htbp}{B05_datenmenge}{Gespeicherte Daten je Lauf und Dateigrößenverteilung der PDFs (logarithmische Achse).}

\section{Wie viele Hochschulen und Studiengänge werden erfasst?}
Von \zUnisTotal{} Hochschulen haben \zUnisStrict{} (\zUnisStrictPct\,\%) mindestens ein LLM-bestätigtes Handbuch; mit den ungeprüften Treffern oberhalb von 0,9 sind es \zUnisLikely{} (\zUnisLikelyPct\,\%). Bei \zUnisGeTen{} Hochschulen liegen mindestens zehn Handbücher vor, bei \zUnisNone{} Hochschulen wurde \emph{nichts} gefunden (kein Dokument, keine Seeds). Die Verteilung ist stark schief (\cref{fig:10_mh_pro_hochschule}): Der Median liegt bei \zMedianMH{} Handbüchern je Hochschule, die zehn ergiebigsten Hochschulen liefern \zTopTenShare\,\% aller Treffer. Ein Host allein (der Modulhandbuch-Server der Universität Augsburg) stellt über 1.800 Dokumente (\cref{fig:17_top_hosts}).

\abb[0.85]{htbp}{10_mh_pro_hochschule}{Verteilung der Treffer pro Hochschule.}
\abb[0.85]{htbp}{17_top_hosts}{Die 20 ergiebigsten Hosts.}

\paragraph{Studiengänge.} Ein Modulhandbuch beschreibt in der Regel einen Studiengang (bzw.\ eine Prüfungsordnungsversion); die Zahl der Handbücher ist daher eine Näherung an die Zahl der erfassten Studiengänge, aber \emph{keine} exakte Zählung: Es gibt Sammelhandbücher für mehrere Studiengänge, mehrere Versionen desselben Studiengangs und Handbücher für Teilmodule. Eine saubere Zählung der Studiengänge wäre erst mit einer inhaltlichen Extraktion möglich (Ausblick, Kapitel~\ref{ch:fazit}).

\section{Wie unterscheiden sich kleine und große, öffentliche und private Hochschulen?}
Die Abdeckung hängt stark von Typ, Träger und Größe der Hochschule ab (\cref{fig:11_abdeckung_typ,fig:12_abdeckung_traeger,fig:13_abdeckung_groesse}):
\begin{itemize}
  \item \textbf{Typ:} Universitäten \zCovUni\,\%, Künstlerische Hochschulen \zCovKunst\,\%, Fachhochschulen/HAW \zCovFH\,\%, Verwaltungshochschulen \zCovVerw\,\%.
  \item \textbf{Träger:} öffentlich-rechtlich \zCovOeff\,\%, kirchlich \zCovKirch\,\%, privat \zCovPriv\,\%.
  \item \textbf{Größe:} Hochschulen unter 1.000 Studierenden liefern im Mittel nur rund \zMHProUniKlein{} Handbücher und ca.\ \zDocsProUniKlein{} heruntergeladene PDFs; Hochschulen mit 15.000--30.000 Studierenden im Mittel rund \zMHProUniGross{} Handbücher bei rund \zDocsProUniGross{} PDFs (\cref{fig:B07_klein_gross}).
\end{itemize}
Der Zusammenhang ist plausibel: Kleine private Hochschulen veröffentlichen Handbücher seltener oder in Login-Bereichen, haben kaum PDFs, nutzen JavaScript-Portale oder gar keine öffentlichen Handbücher. Große öffentliche Hochschulen haben oft Modul-Datenbanken mit tausenden Einzeldokumenten. Das heißt auch: Der Aufwand pro gefundenem Handbuch ist bei kleinen Hochschulen höher, der Nutzen je Crawl-Minute geringer. Grafik~\ref{fig:16_groesse_vs_ausbeute} zeigt die Streuung (Studierendenzahl gegen Treffer).

\abb[0.8]{htbp}{11_abdeckung_typ}{Abdeckung nach Hochschultyp.}
\abb[0.8]{htbp}{12_abdeckung_traeger}{Abdeckung nach Trägerschaft.}
\abb[0.8]{htbp}{13_abdeckung_groesse}{Abdeckung nach Hochschulgröße (Studierendenzahl).}
\abb[0.95]{htbp}{B07_klein_gross}{Mittlere Treffer, mittlere PDFs und Anteil ergiebiger Hochschulen nach Größenklasse.}
\abb[0.8]{htbp}{16_groesse_vs_ausbeute}{Studierendenzahl gegen Treffer pro Hochschule.}
\abb[0.8]{htbp}{14_abdeckung_bundesland}{Abdeckung nach Bundesland.}

\section{Wie hoch ist die Erkennungs- und Klassifikationsquote?}
Die Antwort hängt davon ab, was als Referenz gilt (vgl.\ Kapitel~\ref{ch:klassifikation}):
\begin{itemize}
  \item \textbf{Trainingsmetrik} (gruppierte Kreuzvalidierung): Precision \zTrainPrec\,\%, Recall rund 95\,\%.
  \item \textbf{Realer Crawl, Score 0,5--0,9:} Precision rund \zMidPrec\,\% (Referenz LLM).
  \item \textbf{Realer Crawl, Score $\ge$ 0,9:} Precision \zPilotPct\,\% (Pilot, $n=\zPilotN$).
  \item \textbf{Endliste erweitert:} geschätzt rund \zEstTrue{} echte von \zFinalListed{} gelisteten Dokumenten, das entspricht einer Precision von rund 89\,\%; die streng bestätigte Liste ist per Definition vom LLM bestätigt; ihre tatsächliche Fehlerquote entspricht der Fehlerquote des LLM, die erst die manuelle Stichprobe messen kann (die Referenz ist das LLM, nicht ein Mensch).
  \item \textbf{Recall:} unbekannt, nur grob abschätzbar: Gegenüber dem Erwartungswert von rund \zZielwert{} Handbüchern liegt die geschätzte Zahl echter Treffer bei rund 40--47\,\%. Aus den nicht LLM-geprüften Negativen mit Score 0,3--0,5 (18.252 Dokumente) lässt sich eine kleine Stichprobe von rund \zRecallMissed{} verpassten Handbüchern schätzen, \emph{sehr unsicher} (\cref{fig:21_recall_luecke}).
\end{itemize}

\abb[0.85]{htbp}{21_recall_luecke}{Was vorliegt und was fehlen könnte (grobe Schätzung).}

\section{Wie viele Dokumente müssen manuell nachbearbeitet werden?}
Ohne die LLM-Zweitstufe müssten alle Dokumente der Grauzone manuell geprüft werden. Mit ihr bleiben:
\begin{itemize}
  \item \zUnresolved{} Dokumente, bei denen weder TF-IDF noch OCR/LLM ein Urteil liefern konnten (kein Text, Fehler);
  \item eine Stichprobe von 120 Dokumenten zur Messung der LLM-Genauigkeit (\texttt{stichprobe\_manuelle\_pruefung.csv}, vorbereitet, noch nicht ausgewertet \TODO{Stichprobe auswerten und Ergebnis ergänzen: Übereinstimmung Mensch--LLM});
  \item die 159 verdächtigen Hosts mit rund 6.464 ungeprüften Positiven (LLM-Bestätigungsrate unter 30\,\% bei den übrigen Dokumenten dieser Hosts), deren LLM-Prüfung zum Zeitpunkt des Berichts läuft.
\end{itemize}
Bei einer Annahme von 30 Sekunden je Dokument ersparte die LLM-Zweitstufe rund \zManuellOhneH{} Stunden manueller Prüfung (\cref{fig:B09_manueller_aufwand}; Modellrechnung, keine Messung).

\abb[0.8]{htbp}{B09_manueller_aufwand}{Manueller Prüfaufwand ohne und mit LLM-Zweitstufe (Modellrechnung, 30\,s je Dokument).}

\section{Wie viele Duplikate gibt es?}
In der erweiterten Liste sind \zDupLikely{} Einträge Duplikate (gleicher Host und gleicher Dateiname), in der strengen Liste \zDupStrict{} (\cref{fig:22_duplikate}). Eindeutig bleiben \zUniqueLikely{} (erweitert) bzw.\ \zUniqueStrict{} (streng) Dokumente. Die Duplikate entstehen, weil dieselbe Datei über mehrere Seiten verlinkt ist oder in mehreren Läufen vorkommt; auf Speicherebene kommen die physischen Duplikate erzwungener Läufe hinzu (siehe oben). Die Deduplikation per Host und Dateiname ist eine Näherung; zwei Hochschulen mit gleichnamigen Dateien auf getrennten Hosts bleiben richtig getrennt.

\abb[0.8]{htbp}{22_duplikate}{Duplikate in der Endliste.}

\section{Wie weit reichen die ältesten Dokumente zurück?}
Das Erscheinungsjahr ist nur indirekt bestimmbar. Bei \zJahrAnzahl{} der gelisteten Dokumente steht eine Jahreszahl im Dateinamen. Sie reicht von \zJahrMin{} (nur \zJahrMinAnzahl{} Dokumente) bis 2027 (Handbücher für künftige Semester). Der Median ist \zJahrMedian, \zJahrAnteilNeu\,\% der Dokumente tragen ein Jahr ab 2020, \zJahreMitte{} die Jahre 2015 bis 2019 und \zJahreAlt{} ein Jahr vor 2015 (\cref{fig:19_jahre_im_dateinamen}). Der Crawl ist also zu den aktuellen Jahrgängen verzerrt, hat aber einen langen Schwanz älterer Versionen: Hochschulen lassen alte Handbuchversionen online. Diese Angabe ist ein Näherungswert: Eine Jahreszahl im Namen ist nicht notwendig das Veröffentlichungsjahr, und viele Dateien haben gar keine.

\abb[0.85]{htbp}{19_jahre_im_dateinamen}{Jahreszahlen in den Dateinamen der gelisteten Dokumente.}

\section{In welcher Tiefe der Webseiten liegen die Dokumente?}
Die Klicktiefe wurde im Manifest nicht protokolliert. Als Näherung dient die \emph{Pfadtiefe der Fundseite} (Anzahl der URL-Segmente der Seite, die auf das Dokument verlinkt). Der häufigste Wert über alle PDFs ist \zTiefeModusAlle, bei den gelisteten Handbüchern \zTiefeModusMH; \zTiefeMHBisVier\,\% der Handbücher liegen auf Fundseiten mit Pfadtiefe bis~4, Handbücher liegen also im Mittel \emph{tiefer} als Durchschnitts-PDFs (\cref{fig:B06_tiefe}). Das stützt die Erfahrung aus Kapitel~\ref{ch:crawler}: Eine feste Tiefe von 2, wie sie der erste Crawl nutzte, verpasst den Großteil. \TODO{Für eine echte Klicktiefenstatistik die Tiefe im Manifest mitschreiben (Scrapy-Meta \texttt{depth}); für diesen Bericht Näherung belassen}

\abb[0.85]{htbp}{B06_tiefe}{Pfadtiefe der Fundseiten (Näherung für die Klicktiefe) für alle PDFs und die gelisteten Modulhandbücher.}

\section{Welche Datenmenge wäre bei einer deutschlandweiten Vollerhebung zu erwarten?}
Die Hochschulliste umfasst \zUnisTotal{} Hochschulen und damit praktisch das gesamte Hochschulwesen; die \emph{Breite} der Erhebung ist also bereits deutschlandweit, die \emph{Tiefe} noch nicht. Für den Erwartungswert von rund \zZielwert{} Modulhandbüchern (Fachschätzung) liefert eine lineare Hochrechnung bei gleicher Ausbeute (\zPDFsProMH{} PDFs je echtem Modulhandbuch) rund \zPDFsHochrechnung{} zu sammelnde PDFs, rund \zGBHochrechnung{}\,GB Speicher und rund \zLLMDocsHochrechnung{} LLM-geprüfte Dokumente mit Kosten von etwa \zLLMKostenHochrechnung{}\,\$ (\cref{fig:B08_hochrechnung}). Die Rechenkosten des Crawls bleiben auch dann einstellig. Die Hochrechnung ist eine \emph{Untergrenze}: Die noch fehlenden Handbücher liegen auf den schwer erreichbaren Seiten (kleine, private, JavaScript, Login), die Ausbeute pro PDF wird dort schlechter sein. Wesentlich realistischer als mehr Crawl-Volumen ist daher ein anderer Weg zu den fehlenden Handbüchern (Discovery, Seitenebene, manuelle Nachsuche in der Liste der 39 leeren und der 117 Hochschulen ohne bestätigtes Handbuch).

\abb[0.95]{htbp}{B08_hochrechnung}{Lineare Hochrechnung der Erhebung für \zZielwert{} Modulhandbücher (Ist: durchgezogen, Hochrechnung: schraffiert).}

\chapter{Diskussion: Sackgassen, Lehren und Grenzen}\label{ch:diskussion}

Das Projekt hat mehr Wege verworfen, als es am Ende beschritten hat. Dieses Kapitel sammelt sie bewusst an einer Stelle, ordnet sie ein und zieht Schlüsse. Es folgt eine ehrliche Bewertung der Aussagekraft der Ergebnisse.

\section{Wege, die nicht (oder nur teilweise) funktioniert haben}
\Cref{tab:sackgassen} listet die Sackgassen und teilweisen Fehlschläge in der Reihenfolge ihres Auftretens. Der \emph{Nutzen} der Sackgassen liegt in dem, was sie gezeigt haben: oft die Ursache des eigentlichen Problems.

\begin{table}[htbp]
\centering\small
\caption{Verworfene oder gescheiterte Wege und ihre Lehre.}\label{tab:sackgassen}
\begin{tabularx}{\linewidth}{@{}p{0.5cm} p{3.4cm} p{3.6cm} X@{}}
\toprule
\textbf{Nr.} & \textbf{Weg} & \textbf{Befund} & \textbf{Lehre / Ersatz} \\
\midrule
1 & Lambda für alle Läufe, Tiefe/Seiten hochdrehen & 15-min-Limit, RAM-Pufferung; Timeouts bei Tiefe $\ge$ 2 mit vielen Seiten & Fargate-Bulk-Pfad; Lambda nur noch interaktiv \\
2 & Synchrone Webapp-Antwort & HTTP-Timeout bei längeren Läufen; Endlosschleife bei Mehrfach-Formaten & Asynchrone Jobs mit Polling; Visited-Store; harte Obergrenzen \\
3 & Terraform als Infrastruktur-Wahrheit & Zustand und Handanlagen liefen auseinander; Werkzeuge im Dev-Setup nicht nutzbar & AWS-CLI + JSON-Vorlagen, später Bootstrap-Skript \\
4 & Ungesteuerter Tiefen-Crawl & TU Dortmund: 4 echte von 934 Dokumenten & Best-First, Sitemap, Subdomain-Erkennung \\
5 & Sitemaps als Quelle für tiefe Dateien & Konstanz hat keine Sitemap; tiefe Dateien stehen nirgends & Suchindizes (Common Crawl, Serper) \\
6 & Google Custom Search & Web-weite Suche entfernt (50-Domain-Grenze) & Serper (SERP-API) \\
7 & DuckDuckGo-Scraping & CAPTCHA aus der Cloud & verworfen \\
8 & Firecrawl-Rendering (JS) & Für IU keine PDFs; Pipeline gebaut, nicht im Großlauf getestet & offen; Playwright-basiertes Rendering als Ausblick \\
9 & Schwelle 0,5 für maximalen Recall & Positiv-Flut, Precision 25--60\,\% in der Praxis & LLM-Zweitstufe; Schwelle 0,8--0,9 empfohlen \\
10 & OCR für alle Scans & 3.553 Scans, nur 42 bestätigte Handbücher (rund 1\,\%) & Aufwand/Ertrag schlecht; nur für Review-Band sinnvoll \\
11 & Re-Scoring in der Evaluation überschrieb LLM-Urteile & Falsche, zu optimistische Abdeckung (314 statt 303) & LLM-Urteil hat Vorrang; Test-Disziplin \\
12 & DynamoDB-Dedup auf Hochschulebene bei Discovery-Wiederholung & Discovery lief nie für bereits gecrawlte Hochschulen & Seeds-only-Modus \\
13 & Mantle-Endpunkt von Bedrock & per SCP verboten (403) & nativer Bedrock-Client \\
14 & Langlebiger Bedrock-Schlüssel & vom Account nicht erlaubt; Schlüssel verfallen nach 30\,min bis 12\,h & Chunk-Verarbeitung, wiederaufnehmbar \\
15 & Serper bezahlt / \texttt{filetype:} & Free-Tier erlaubt weder \texttt{filetype:} noch mehr als 10 Treffer & Free-Tier-Anpassung; Recall begrenzt \\
16 & SQS-Fan-out für Breite & Roadmap-Punkt, nie gebaut & durch Fargate ersetzt \\
17 & Menschen-Review als Hauptmittel der Qualitätssicherung & Skala (zehntausende Kandidaten) & LLM-Zweitstufe, Mensch nur für Stichprobe \\
\bottomrule
\end{tabularx}
\end{table}

\section{Lehren}
\begin{enumerate}
  \item \textbf{Der Engpass wandert.} Zuerst war es die Plattform, dann die Laufzeit, dann der Recall, am Ende die Precision. Jede Verbesserung deckt das nächste Problem auf. Eine frühe Messung am Ziel (statt auf dem Papier) hätte die Reihenfolge verkürzt.
  \item \textbf{Trainingsmetriken sind keine Betriebsmetriken.} 96,5\,\% Precision in der gruppierten Kreuzvalidierung, 25\,\% im Crawl. Der Unterschied liegt in der Verteilung der Negativbeispiele, nicht im Modell. Wichtigste Lehre für ML-Projekte: eine Stichprobe aus der \emph{echten} Verteilung früh labeln.
  \item \textbf{Eine zweite, stärkere Stufe lohnt für die Grauzone.} Ein billiges Modell filtert, ein teures prüft nur die Unsicheren. Das spart gut 80\,\% der LLM-Aufrufe gegenüber einer Vollprüfung.
  \item \textbf{Die Evaluation braucht eigene Tests.} Der Fehler, der LLM-Urteile überschrieb, war nicht im Modell, sondern in der Auswertung, und er verfälschte alle Zahlen zuvor.
  \item \textbf{Erst messen, dann bauen (und danach wieder messen).} OCR und Firecrawl waren plausibel, brachten aber kaum Ertrag; Seeds-only brachte \emph{in Minuten} +43 Hochschulen.
  \item \textbf{Infrastruktur wegwerfbar, Daten nicht.} Der Sandbox-Verlust war verkraftbar, weil der Aufbau skriptbar war; verloren schmerzte, was nur in S3 lag.
  \item \textbf{Randbedingungen sind Anforderungen.} SCPs, kurzlebige Schlüssel und Proxy waren keine Störungen, sondern gaben die Architektur der Nachbereitung vor.
\end{enumerate}

\section{Grenzen und Bedrohungen der Validität}\label{sec:grenzen}
\paragraph{Das LLM ist die Referenz, nicht ein Mensch.} Alle Precision-Angaben messen die Übereinstimmung mit Claude Haiku~4.5. Ein systematischer LLM-Fehler (z.\,B.\ zu strenge Auslegung von „Sammelhandbuch“) steckt in allen Zahlen. Die manuelle Stichprobe von 120 Dokumenten ist vorbereitet und würde die einzige echte Genauigkeitsmessung liefern.

\paragraph{Nicht alles ist LLM-geprüft.} Die \zTfidfHigh{} Positive mit Score $\ge$ 0,9 und die \zTfidfOther{} Positive aus anderen Läufen sind ungeprüft. Die Hochrechnung von 82,5\,\% aus 400 Dokumenten setzt Repräsentativität voraus; 159 Hosts mit zusammen rund 6.464 ungeprüften Positiven haben bei ihren übrigen Dokumenten unter 30\,\% Bestätigungsrate (z.\,B.\ uni-hamburg.de, uni-flensburg.de, tu-braunschweig.de). Dort ist der Anteil Fehlalarme vermutlich höher als der Durchschnitt. Das Ergebnis der laufenden Nachprüfung dieser Hosts verschiebt die Zahlen \TODO{nach Abschluss der LLM-Prüfung Zahlen und Aussage in Abschnitt ergänzen}.

\paragraph{Der Recall ist nur grob bekannt.} Der Erwartungswert von rund \zZielwert{} ist eine Fachschätzung, keine Messung. Es gibt keine unabhängig gezählte Grundgesamtheit pro Hochschule. Die Schätzung verpasster Handbücher unter den Negativen stammt aus einer kleinen, nicht repräsentativen Stichprobe.

\paragraph{Definition.} Ob Sammelhandbücher, Modulkataloge oder Studienverlaufspläne zählen, ist Auslegung. Die strenge Definition (19\,\% der Fehlalarme sind Einzelblätter) ist eine Entscheidung zugunsten der Precision; ein breiterer Begriff würde die Ergebnisliste vergrößern.

\paragraph{Verzerrungen.} (1)~Das Trainingsset ist verzerrt zu Hochschulen mit großem Online-Angebot. (2)~Die Hochschulliste enthält mehr Fachhochschulen und private Hochschulen als Universitäten, die Abdeckung dort ist am schlechtesten. (3)~Der Crawl ist zu aktuellen Jahrgängen verzerrt. (4)~Kosten sind Schätzungen, Stunden sind Schätzungen, die Klicktiefe ist eine Näherung.

\paragraph{Ethik und Recht.} Der Crawler beachtet robots.txt, drosselt Anfragen und greift ausschließlich öffentlich zugängliche Dokumente ab. Die Ausnahme für Discovery-Seeds (robots-Bypass für einzelne, per Suchindex bekannte Handbuch-URLs, Abschnitt~\ref{sec:crawler-robots}) ist standardmäßig aus und wurde nur für Hochschulen mit pauschaler Sperre eingesetzt; sie ist eine bewusste Abwägung, keine Selbstverständlichkeit. Es werden keine personenbezogenen Daten gezielt erhoben. Die Urheberrechtslage bei der Weitergabe der gesammelten Dokumente ist nicht geprüft \TODO{rechtliche Einordnung mit dem Betreuer klären; der Bericht ist kein juristisches Gutachten}.

\chapter{Beitrag von Sascha Lauk}\label{ch:sascha}

\saschaPH{Dieses Kapitel schreibt Sascha Lauk selbst. Die folgende Gliederung ist nur ein Vorschlag, damit Inhalt und Umfang zum restlichen Bericht passen; sie kann frei geändert werden.}

\section{Aufgabenbereich und Einordnung}
\saschaPH{Welche Teile des Systems hat Sascha bearbeitet (z.\,B.\ Kapitelbezug zu Architektur, Crawler, Klassifikation, Evaluation, Frontend)? Wie ist die Schnittstelle zu den Arbeiten in den Kapiteln~\ref{ch:architektur} bis~\ref{ch:klassifikation}?}

\section{Vorgehen und Entscheidungen}
\saschaPH{Eingeschlagene Wege, Alternativen, Sackgassen, wie in Kapitel~\ref{ch:diskussion} für den Hauptteil.}

\section{Ergebnisse und Kennzahlen}
\saschaPH{Messbare Ergebnisse des eigenen Teils (Zahlen, Diagramme, Tabellen).}

\section{Probleme und Lehren}
\saschaPH{Probleme im Betrieb oder in der Entwicklung, Lehren.}

\section{Wochenübersicht}
Die wochenweise Aufschlüsselung von Sascha Lauk steht in Tabelle~\ref{tab:wochen-sascha}.

\chapter{Fazit und Ausblick}\label{ch:fazit}

\section{Fazit}
Das Projekt hat aus einem Scrapy-Prototyp eine skalierbare Erhebungs- und Klassifikationsplattform auf AWS gemacht und sie an einem realen, großen Anwendungsfall erprobt: \zDocsTotal{} PDFs von \zUnisTotal{} Hochschulen, zwei Großläufe mit zusammen rund 18~Stunden Laufzeit, \zSpeicherGB{}\,GB Daten, eine Endliste von \zFinalListed{} Dokumenten mit geschätzt \zEstTrue{} echten Modulhandbüchern und einer Abdeckung von \zUnisLikelyPct\,\% der Hochschulen. Der Aufbau ist modular genug, dass ein anderer Dokumenttyp ein neues Profil und ein neues Modell braucht, nicht aber einen neuen Crawler. Das ist die Brücke zu KIWIT.

\Cref{tab:akzeptanz} bewertet die Akzeptanzkriterien aus Kapitel~\ref{ch:einleitung} gegen das Ergebnis.

\begin{table}[htbp]
\centering\small
\caption{Akzeptanzkriterien und Ergebnis.}\label{tab:akzeptanz}
\begin{tabularx}{\linewidth}{@{}p{2.8cm} p{3.0cm} X@{}}
\toprule
\textbf{Kriterium} & \textbf{Erfüllt?} & \textbf{Begründung} \\
\midrule
A1 vollständig durchsuchen & weitgehend & Alle \zUnisTotal{} Hochschulen wurden bearbeitet; Seiten mit JavaScript, Login oder anderer Domain nicht erreichbar. \\
A2 alle Handbücher finden (Recall) & \textbf{teilweise} & Etwa 40--47\,\% des Erwartungswerts; 39 Hochschulen ohne Fund, 117 mit Dokumenten, aber ohne bestätigtes Handbuch. \\
A3 Nicht-Handbücher erkennen (Precision) & ja, durch LLM-Stufe & TF-IDF allein reicht nicht (25--82\,\%); die streng bestätigte Liste ist vom LLM geprüft, der Rest geschätzt $\approx$ 89\,\%. \\
A4 strukturiert speichern & ja & S3-Ablage mit Manifest (Score, Quelle, Lauf, LLM-Begründung); Transparenz je Entscheidung. Eine Datenbank für Inhalte ist offen. \\
\bottomrule
\end{tabularx}
\end{table}

Das wichtigste inhaltliche Ergebnis ist \emph{negativ} und nützlich: Der Klassifikator, so gut er auf dem Papier war, genügte als alleiniger Filter nicht. Die Kombination aus billigem Massenfilter und LLM-Zweitstufe war die Antwort, und die Ehrlichkeit der Evaluation (gruppiert nach Hochschule, Pilotstichprobe, korrigierter Auswertungsfehler) hat dieses Ergebnis überhaupt erst sichtbar gemacht. Das wichtigste betriebliche Ergebnis: Eine befristete Sandbox verlangt, dass Aufbau und Daten wiederherstellbar sind.

\section{Ausblick}
\subsection{Qualität und Recall}
\begin{itemize}
  \item \textbf{Schwelle und Modell.} Im Live-Crawl die vorgeschlagene Schwelle (0,8) statt 0,5 verwenden; den Dateinamen als Merkmal aufnehmen (Abschnitt~\ref{ch:klassifikation}); mit den rund 27.700 LLM-Labels als \emph{Hard Negatives} nachtrainieren. Die dazu nötige Retrain-Strecke aus Manifesten (\texttt{ingest --manifest} mit S3-Abruf) ist vorhanden.
  \item \textbf{Manuelle Referenz.} Die vorbereitete Stichprobe von 120 Dokumenten auswerten, um die LLM-Genauigkeit zu messen.
  \item \textbf{Verdächtige Hosts.} Die ungeprüften Positive der 159 Hosts mit niedriger Bestätigungsrate LLM-prüfen (rund 15--20\,\$).
  \item \textbf{Seitenebene statt Dateiebene.} JavaScript-Portale (z.\,B.\ IU) liefern kein PDF; die Klassifikation von gerenderten Katalogseiten ist gebaut und wartet auf einen Einsatz mit Playwright oder Firecrawl.
  \item \textbf{Manuelle Nachsuche} für die 39 leeren und die kleinen privaten Hochschulen, deren Handbücher nicht online liegen oder hinter Login.
\end{itemize}

\subsection{Inhaltliche Erschließung}
Aus „ist ein Modulhandbuch“ soll „was steht darin“ werden: Module, ECTS-Punkte, Lernziele und Prüfungsformen als strukturierte, vergleichbare Daten. Für diese Extraktion ist ein LLM fachlich naheliegend; erst damit lassen sich Studiengänge zählen und systematisch vergleichen, wie es der ursprüngliche Anwendungsfall der FH Wedel verlangt.

\subsection{Plattform und Betrieb}
\begin{itemize}
  \item Modell-Artefakte in S3 mit Registry versionieren, sodass Modellwechsel ohne Neubau des Images möglich ist.
  \item Scoring vom Crawl entkoppeln (S3-Ereignis lösen Klassifikations-Lambda aus) und Jobstatus in einer Tabelle halten (Roadmap Phase~3--4).
  \item Secrets Manager statt \texttt{.env}-Dateien, echtes TLS-Zertifikat statt selbstsigniertem, SQLite $\to$ RDS bei mehreren Schreibern, Auto-Scaling hinter einem Load Balancer.
  \item Automatisierte Tests und Monitoring (Fehlerraten, Review-Band-Anteil als Drift-Indikator).
  \item Webapp nur bei Bedarf betreiben: Die dauerhaft laufende Instanz kostete im Budget mehr als jeder Crawl; ein On-Demand-Betrieb wäre günstiger.
\end{itemize}

\subsection{Anschluss an KIWIT}
Für Forschungsprofil~A der KIWIT-Gruppe (Makroebene) lassen sich damit künftig auch KI-bezogene Strategien, Richtlinien und Angebote erheben: Dazu genügen ein neues Extraktionsprofil mit Suchbegriffen, eine neue Spezifikation für die Zweitstufe und ein passendes Trainingsset. Die in diesem Projekt gewonnene Erfahrung, dass \emph{Auffinden} und \emph{Aussortieren} getrennte Probleme sind und beide eine eigene Lösung brauchen, sollte dabei von Beginn an berücksichtigt werden.


\appendix
\chapter{Konfiguration der Läufe}\label{app:konfig}

\begin{table}[htbp]
\centering\small
\caption{Konfigurationsgrenzen des Crawlers: Lambda/MVP gegenüber Fargate Deep Run.}
\begin{tabularx}{\linewidth}{@{}l l l X@{}}
\toprule
\textbf{Parameter} & \textbf{MVP/Lambda} & \textbf{Deep Run} & \textbf{Bedeutung} \\
\midrule
\texttt{CRAWL\_MAX\_DEPTH} & 2 & 5 & Link-Tiefe je Hochschule \\
\texttt{CRAWL\_MAX\_PAGES} & 60 & 2.500 & Seitenbudget (weich) \\
\texttt{CRAWL\_MAX\_PAGES\_PER\_SECTION} & -- & 60 & Seitenkappe je Fakultäts-Sektion \\
\texttt{CLOSESPIDER\_PAGECOUNT} & 400 & 3.000 & harte Seitengrenze \\
\texttt{CLOSESPIDER\_TIMEOUT} & -- & 1.800\,s & Zeitgrenze je Hochschule \\
\texttt{BULK\_MAX\_RUNTIME\_SECONDS} & -- & 43.200 & Wachhund (Gesamtlauf) \\
\texttt{CRAWL\_MAX\_SUBDOMAIN\_SITEMAPS} & 15 & 30 & Sitemaps von Subdomains \\
\texttt{DOWNLOAD\_MAXSIZE} & 50\,MB & 50\,MB & je Datei \\
\texttt{DOWNLOAD\_TIMEOUT} & 60\,s & 60\,s & je Anfrage \\
\texttt{DOWNLOAD\_DELAY} & 1\,s + AutoThrottle & dto. & je Domain \\
\texttt{CONCURRENT\_REQUESTS} & 4 (2 je Domain) & 16 & je Job \\
Rechenressource & Lambda 1.536\,MB & Fargate 4\,vCPU/16\,GB & \\
\texttt{CLASSIFIER\_UPPER\_THRESHOLD} & 0,65 & 0,50 & obere Modellschwelle \\
\bottomrule
\end{tabularx}
\end{table}

\chapter{Metriken im Detail}\label{app:metriken}
\begin{table}[htbp]
\centering\small
\caption{Metriken der Modellvarianten (gruppierte Kreuzvalidierung), jeweils am Recall-Ziel von 95\,\%.}
\begin{tabular}{@{}l r r r r r r@{}}
\toprule
\textbf{Variante} & \textbf{PR-AUC} & \textbf{ROC-AUC} & \textbf{Recall} & \textbf{Prec.} & \textbf{F1} & \textbf{F2} \\
\midrule
MVP \emph{content} (0,387) & 0,957 & 0,941 & 0,955 & 0,871 & 0,911 & 0,937 \\
MVP \emph{content + Dateiname} & 0,960 & 0,926 & 0,955 & 0,848 & 0,898 & 0,932 \\
Sept.\ \emph{content} (0,627) & 0,988 & 0,973 & 0,950 & 0,965 & 0,958 & 0,953 \\
Sept.\ \emph{content + Metadaten} (0,509) & 0,976 & 0,959 & 0,950 & 0,958 & 0,954 & 0,952 \\
\bottomrule
\end{tabular}
\end{table}
\paragraph{Konfusionsmatrizen.} MVP \emph{content}: 128 richtig positiv, 6 falsch negativ, 49 richtig negativ, 19 falsch positiv (202 Dokumente, 21 Hochschulen); bei Schwelle 0,5: Precision 0,907, Recall 0,873, 17 statt 6 übersehene Handbücher. September \emph{content}: 1.873 / 98 / 656 / 68 (2.695 Dokumente, 269 Hochschulen); bei Schwelle 0,5: Precision 0,952, Recall 0,985 (1.942 TP, 29 FN, 625 TN, 99 FP).

\chapter{Reproduzierbarkeit}\label{app:repro}
Alle Zahlen und Grafiken lassen sich aus den Läufen neu erzeugen:
\begin{lstlisting}
# 1. Endliste und Kennzahlen (aus Repo-Root)
python auswertung/build_final.py
.venv\Scripts\python auswertung\make_charts.py       # 23 Grafiken + kennzahlen.json

# 2. Berichtsgrafiken und Zahlenmakros
.venv\Scripts\python bericht\make_report_charts.py   # B01-B09, zahlen.tex, Wochentabellen

# 3. PDF bauen
cd bericht ; latexmk -pdf main.tex
\end{lstlisting}
Die gesamte Strecke steht in \texttt{bericht/build.ps1}. Handgepflegte Eingaben: \texttt{bericht/daten/stunden.csv} (Stunden und Tätigkeiten), \texttt{bericht/daten/kosten.csv} (Kosten), \texttt{bericht/daten/s3\_objekte.csv} (Auflistung des Buckets, erzeugt mit \texttt{aws s3 ls --recursive}). Evaluationsberichte (\texttt{eval\_*.html}) entstehen mit \texttt{python -m webscraper.evaluation} aus \texttt{webscraper/}.

\chapter{Weitere Auswertungsgrafiken}\label{app:grafiken}
Die folgenden Grafiken der Abschlussauswertung werden im Text nicht direkt referenziert.
\abb[0.85]{htbp}{04_score_histogramm}{Score-Verteilung nach Endstatus (logarithmische Achse).}
\abb[0.85]{htbp}{06_llm_konfidenz}{Konfidenzverteilung der LLM-Urteile.}
\abb[0.85]{htbp}{15_hochschul_status}{Status der Hochschulen vor und nach der Bereinigung.}

\chapter{Glossar}
\begin{description}[style=nextline,leftmargin=1.2cm]
  \item[Bulk-Crawl] Läufe über viele Hochschulen auf Fargate ohne Zeitlimit.
  \item[Discovery] Vorabsuche nach Handbuch-URLs über Suchindizes (Common Crawl, Serper).
  \item[Seeds] Von der Discovery gefundene URLs, die der Crawler priorisiert holt.
  \item[Seeds-only-Modus] Lauf, der ausschließlich Seeds holt und den Visited-Store umgeht.
  \item[Visited-Store] Persistenter Speicher (DynamoDB) für bereits gecrawlte Hochschulen.
  \item[Review-Band] Score-Bereich zwischen den Schwellen; Dokumente mit unsicherer Entscheidung.
  \item[Pilotstichprobe] Zufällige Stichprobe von 400 Positiven mit Score $\ge$ 0,9 zur Schätzung der Precision der ungeprüften Gruppe.
  \item[Strenge / erweiterte Liste] Strenge Liste: nur LLM-bestätigt. Erweiterte Liste: zusätzlich TF-IDF $\ge$ 0,9 ohne LLM-Prüfung.
  \item[SCP] Service Control Policy: organisationsweite AWS-Sperre einzelner Aktionen.
  \item[KIWIT] Forschungsgruppe zu Funktionen und Folgen Künstlicher Intelligenz in der Wissenschafts- und Hochschulorganisation.
\end{description}


\bibliographystyle{plainnat}
\bibliography{literatur}

\end{document}
